\chapter{Partial Okounkov bodies and Duistermaat--Heckman measures}
\index{Okounkov body!partial}\index{Duistermaat--Heckman measure}\index{numerical dimension!toric}
\label{chap:transcendental-okounkov-dh}

Okounkov bodies turn asymptotic spaces of sections into convex sets.
For singular metrics the relevant spaces are cut out by multiplier ideals,
and for transcendental classes they must be recovered by approximation.
The resulting body remembers the multiplier-ideal singularity type and its
Euclidean volume records the volume detected by multiplier ideals.  This is
not, for an arbitrary singular potential, the same as its raw
non-pluripolar mass; the distinction from \Cref{chap:complex-prerequisites}
is essential throughout this chapter.
Applied level by level to a test curve, these bodies carry a concave test
function whose distribution is the Duistermaat--Heckman measure.

The algebraic construction follows the published paper \cite{Xia21}.
The transcendental construction and the positive-mass DH package use
\cite[Sections~10.4--10.5 and~13.3]{XiaBook}; the ordinary transcendental
body theorem is the interface summarized in its Theorem~10.4.1 from
\cite{DRWNXZ}.  Chapters~51--54 supply the envelope, mass, test-curve
and trace conventions used below.
The mixed-volume section, \Cref{sec:ok-mixed-volumes}, proves
the positive-volume formula conjectured in \cite[Conjecture~8.8]{Xia21}.
The proof begins with the argument in the work in progress \cite{XiaMVOB},
repairs its convex-sandwich and approximation steps, and records every
external input at the point where it is used.  Its scope is projective,
algebraic, and strictly positive-volume; no zero-volume or transcendental
extension is inferred from it.
The last section has a different status.
It records a validated current-level toric theorem in arbitrary real classes,
an independent section-counting proof in rational classes, and one obstruction
from the research ledger.  None is obtained from the positive-volume partial
body theory by formal continuity.

Throughout, $X$ is a connected compact K\"ahler manifold of dimension
$n\geq1$.
In algebraic statements $X$ is smooth and projective.
Lebesgue measure on $\mathbf R^n$ is denoted by $\lambda_n$.

\section{The algebraic partial body}
\label{sec:ok-algebraic-body}

Let $L$ be a big line bundle on a smooth projective variety $X$.
Choose a smooth Hermitian metric $h_0$ with curvature $\theta$, and write
\[
                         h=h_0\rme^{-\varphi},
        \qquad \varphi\in\PSH(X,\theta).
\tag{55.1}\label{eq:ok-partial-metric}
\]
Fix a surjective rank-$n$ valuation
\[
             \nu:\mathbf C(X)^\times\longrightarrow\mathbf Z^n_{\mathrm{lex}}
\]
trivial on $\mathbf C^\times$, with one-dimensional leaves.
Surjectivity fixes the full valuation lattice and excludes an extra
lattice-index factor in the volume formula.
An admissible flag gives such a valuation by successive orders of vanishing.
For a general valuation, choose a rational frame $e$ of $L$ and put
$\nu(s)=\nu(s/e^k)$ for a section of $L^k$.
Changing $e$ translates every normalized body by one fixed vector,
as proved in \Cref{lem:ok-mvob-valuative-leaves}.

\begin{definition}[Multiplier-ideal series]
\label{def:ok-multiplier-series}
For $k\geq1$, put
\[
 W_k(\theta,\varphi)
 :=H^0\!\left(X,L^k\otimes\cI(k\varphi)\right).
\]
The valued graded family attached to $(\theta,\varphi)$ is
\[
 \Gamma_\nu(\theta,\varphi)
 :=\{(\nu(s),k):k\geq1,\ s\in W_k(\theta,\varphi)\setminus\{0\}\}.
\tag{55.2}\label{eq:ok-valued-family}
\]
\end{definition}

Multiplier ideals are only submultiplicative in the direction
\[
 \cI((k+\ell)\varphi)
       \subseteq\cI(k\varphi)\cI(\ell\varphi).
\]
Thus the product of sections at levels $k$ and $\ell$ need not belong
to level $k+\ell$.  Consequently $\Gamma_\nu(\theta,\varphi)$ need not
be an honest semigroup with all the usual saturation properties.
The almost-semigroup construction nevertheless supplies a body.  Here is
its exact imported interface, \cite[Theorem~3.14 and Corollary~3.15]{Xia21},
or \cite[Theorem~10.1.2 and Corollary~10.1.2]{XiaBook}.
Work in a fixed closed pointed cone with compact slice at degree one.
For two valued families $A,B$ in that cone, set
\[
 d_{\mathrm{sg}}(A,B)
 :=\limsup_{k\to\infty}k^{-n}\#(A_k\mathbin{\triangle}B_k),
 \qquad A_k:=\{a:(a,k)\in A\}.
\tag{55.2a}\label{eq:ok-almost-semigroup-distance}
\]
An almost semigroup is a positive-counting-volume limit in this
pseudometric of semigroups generating $\mathbf Z^{n+1}$.
The ordinary cone-slice body map on those semigroups has a unique
continuous extension to this positive-volume completion.
The extension preserves counting volume and inclusion.
The multiplier-ideal family above belongs to that completion by
\cite[Section~5.1]{Xia21}, or \cite[Theorem~10.3.1]{XiaBook}.
Thus the symbol $\Delta(\Gamma_\nu)$ below denotes this asymptotic
extension, rather than the cone generated by all its finite-degree points.

\begin{definition}[Algebraic partial Okounkov body]
\label{def:ok-algebraic-partial-body}
Assume $\vol(\theta,\varphi)>0$, in the sense of
\Cref{def:complex-current-volume}.
The partial Okounkov body is
\[
 \Delta_\nu(\theta,\varphi)
       :=\Delta\!\left(\Gamma_\nu(\theta,\varphi)\right)
       \subset\mathbf R^n.
\tag{55.3}\label{eq:ok-partial-body}
\]
For a flag $Y_\bullet$ we write $\Delta_{Y_\bullet}$.
\end{definition}

The positive-volume assumption is structural here.
It makes the body full-dimensional and is used by the continuity theorem
for almost semigroups.
The definition does not silently assign a full-dimensional body when the
multiplier-ideal volume is zero.

\begin{theorem}\label{thm:ok-algebraic-volume-order}
Let $\varphi,\psi\in\PSH(X,\theta)$ have positive multiplier-ideal
volume.
Then
\[
 \lambda_n\!\left(\Delta_\nu(\theta,\varphi)\right)
     =\frac1{n!}\vol(\theta,\varphi).
\tag{55.4}\label{eq:ok-algebraic-volume}
\]
If
$\varphi\preceq_{\cI}\psi$, meaning
$\cI(k\varphi)\subseteq\cI(k\psi)$ for every $k>0$, then
\[
 \Delta_\nu(\theta,\varphi)
       \subseteq\Delta_\nu(\theta,\psi).
\tag{55.5}\label{eq:ok-i-order}
\]
In particular the body depends only on the $\cI$-singularity type.
For a potential with minimal singularities it is the usual Okounkov body
$\Delta_\nu(L)$.
\end{theorem}

\begin{proof}[Almost-semigroup source boundary]
The ideal inclusion gives
$W_k(\theta,\varphi)\subseteq W_k(\theta,\psi)$ for every $k$.
Monotonicity of the almost-semigroup body gives
\eqref{eq:ok-i-order}.

The remaining assertion is the central almost-semigroup theorem of
\cite[Section~5.1]{Xia21}, with the precise almost-semigroup interface
recorded above and \cite[Corollary~10.3.2]{XiaBook}.
We record its inputs so that the boundary is visible.
First, the multiplier-ideal family
$\Gamma_\nu(\theta,\varphi)$ is proved to be an almost semigroup, and its
asymptotic counting volume is
\[
 \lim_{k\to\infty}k^{-n}
 h^0\!\left(X,L^k\otimes\mathcal I(k\varphi)\right)
 =\frac1{n!}\vol(\theta,\varphi).
\]
This is \cite[Theorem~2.19]{Xia21}, in the normalization of
\eqref{eq:complex-current-volume}.  The one-dimensional leaves identify
section dimensions with valuation counts.  Second, the
almost-semigroup body construction is continuous for its pseudometric and
has Euclidean volume equal to the counting volume.

For orientation, the source proves these inputs first for analytic
singularities defining a K\"ahler current.  On a log resolution
\[
                   \pi^*\theta_\varphi=[D]+R,
\tag{55.6}\label{eq:ok-log-resolution}
\]
divisorial valuations translate the ordinary Okounkov body by
$\nu(D)$, and asymptotic Riemann--Roch computes the residual volume.  It
then passes to K\"ahler currents by quasi-equisingular approximation and
to arbitrary positive-volume potentials by
$\varphi_\varepsilon=(1-\varepsilon)\varphi+\varepsilon\chi$.
Continuity is required both for the almost semigroups and for
$\mathcal I$-model volume.  Applying the two imported inputs to
\eqref{eq:ok-partial-body} gives
\eqref{eq:ok-algebraic-volume}.
\end{proof}

\begin{example}[A pole on the projective line]
\label{ex:ok-p1-partial-interval}
On $\mathbf P^1$, take $L=\mathcal O(1)$ and a section $s_p$ vanishing
at the flag point $p$.  For $0<a<1$, set
$\varphi=a\log|s_p|_{h_0}^2$, with $c_1(L,h_0)=\omega_{\mathrm{FS}}$.
Then
\[
 \theta_\varphi=(1-a)\omega_{\mathrm{FS}}+a\delta_p,
 \qquad \mathcal I(k\varphi)=\mathcal O(-\lfloor ka\rfloor p).
\tag{55.6a}\label{eq:ok-p1-pole-ideals}
\]
The ideal formula follows by integrating a monomial $z^j$ against
$|z|^{-2ka}$: local integrability is $j>ka-1$.
Thus the valuation set at degree $k$ is
$\{\lfloor ka\rfloor,\ldots,k\}$.
Replacing its lower bound by $\lceil ka\rceil$ gives an honest semigroup:
$\lceil(k+\ell)a\rceil\leq\lceil ka\rceil+\lceil\ell a\rceil$.
It generates $\mathbf Z^2$, since it contains $(1,1)$ and consecutive
values at sufficiently large degrees.  Its closed cone slice is $[a,1]$.
At most one point per degree has been changed, so the distance in
\eqref{eq:ok-almost-semigroup-distance} is zero.  Consequently
\[
 \Delta_p(\theta,\varphi)=[a,1],\qquad
 \lambda_1(\Delta_p)=1-a=M_\theta(\varphi)=\vol(\theta,\varphi).
\tag{55.6b}\label{eq:ok-p1-partial-body}
\]
The raw mass removes the atom at $p$, and its equality with ideal volume
uses analytic singularities.  At degree one the original valuation set
contains zero.  Taking the cone of all finite-degree points would therefore
produce $[0,1]$ instead, which explains the asymptotic definition.
\end{example}

\begin{remark}[Why raw mass cannot occur in \eqref{eq:ok-algebraic-volume}]
\label{rem:ok-ideal-volume-necessity}
Every graded space in \eqref{eq:ok-valued-family} depends only on the
multiplier ideals of $\varphi$.  Hence the body is constant on an
$\mathcal I$-equivalence class.  Raw non-pluripolar mass is not constant
on such a class: the Cantor example of
\Cref{prop:ok-i-invariant-obstruction} has
$\varphi\sim_{\mathcal I}0$ but
$M_\theta(\varphi)=0<M_\theta(0)$.  Thus replacing
$\vol(\theta,\varphi)$ in \eqref{eq:ok-algebraic-volume} by
$M_\theta(\varphi)$ would contradict the definition of the body.
\end{remark}

\begin{remark}[A closure that cannot be deleted formally]
\label{rem:ok-valuative-closure}
There is a valuative description of the partial body by valuation vectors of
positive currents whose singularity type is dominated by
$\varphi$.
For the transcendental body constructed below, the exact interface is
\cite[Theorem~10.4.4]{XiaBook}: it uses positive currents in the same class,
ordered by $\preceq_{\mathcal I}$, and takes the closure of their
current-valuation vectors.  Current valuations are defined by successive
generic Lelong numbers and traces of the residual currents along the
smooth flag, \cite[Definition~10.2.3]{XiaBook}.
No argument in this chapter removes the closure, and we will never use the
stronger equality without closure.
\end{remark}

\section{The transcendental construction}
\label{sec:ok-transcendental-body}

Let $\alpha\in H^{1,1}(X,\mathbf R)$ be big, let
$\theta\in\alpha$ be smooth, and fix a smooth flag
\[
 X=Y_0\supset Y_1\supset\cdots\supset Y_n=\{p\}.
\]
For a positive current $T=\theta+\ddc\varphi$ of positive
multiplier-ideal volume, write
\[
                       \vol(T):=\vol(\theta,\varphi).
\]
This is independent of the smooth representative $\theta$ by
\Cref{def:complex-current-volume}.  The body
$\Delta_{Y_\bullet}(T)$ is defined in three stages.

\paragraph{Analytic singularities.}
First assume that $T$ is a K\"ahler current with rational analytic
singularities in the sense of \Cref{def:ok-rational-analytic}.
Choose a modification $\pi:Y\to X$ resolving the singularities and lifting
the flag.
Write
\[
                         \pi^*T=[D]+R.
\]
If $Z_\bullet$ is the lifted flag and $g\in\mathrm{GL}_n(\mathbf Z)$
is its change-of-valuation matrix, set
\[
 \Delta_{Y_\bullet}(T)
 :=\bigl(\Delta_{Z_\bullet}(\{R\})+\nu_{Z_\bullet}(D)\bigr)g^{-1}.
\tag{55.7}\label{eq:ok-trans-analytic}
\]
Here $D$ is an effective rational divisor and $R$ is positive with
bounded potentials, so its class is nef and big.
The matrix $g$ is upper triangular with diagonal entries one; valuation
vectors are row vectors, so $|\det g|=1$.
A common resolution and the correction-matrix cocycle prove independence
of choices, \cite[Definition~10.4.2 and Lemma~10.4.1]{XiaBook}.
The ordinary body of $\{R\}$ uses
\cite[Theorem~10.4.1]{XiaBook}, not the multiplier-ideal series.

\paragraph{K\"ahler currents.}
Choose a decreasing quasi-equisingular approximation
$T_j=\theta+\ddc\varphi_j$ inside the same class, with rational analytic
singularities and a common strictly positive curvature lower bound.
If $T\geq\delta\omega$, regularize its potential in background
$\theta-\delta\omega/2$.  The errors tend to zero, so the positive
margin absorbs them and leaves a tail of same-class K\"ahler currents.
This uses the interface of \cite[Theorem~1.6.2]{XiaBook}.
Define
\[
                     \Delta_{Y_\bullet}(T)
                     :=\bigcap_{j\geq1}\Delta_{Y_\bullet}(T_j).
\tag{55.8}\label{eq:ok-trans-kahler}
\]
The independence proof is \cite[Lemma~10.4.4]{XiaBook}.
Its cross-comparison uses $\psi_k\preceq(1-\varepsilon)\varphi_j$
for a sufficiently large $k$, not literal cofinality of the unscaled ideals.
The analytic-body order theorem first gives inclusion in the body of
$(1-\varepsilon)\varphi_j$.  Letting $\varepsilon\downarrow0$ gives
that of $\varphi_j$ by nested-body volume continuity.  Interchanging the
two approximations proves equality of their intersections.

\paragraph{General positive currents.}
Fix a K\"ahler form $\omega$.
For a current $T$ of positive multiplier-ideal volume set
\[
 \Delta_{Y_\bullet}(T)
 :=\bigcap_{j\geq1}
      \Delta_{Y_\bullet}(T+j^{-1}\omega).
\tag{55.9}\label{eq:ok-trans-general}
\]
Each $T+j^{-1}\omega$ is a K\"ahler current, so the preceding stage
applies.  The bodies decrease with $j$.
Changing $\omega$ gives mutually cofinal perturbations by compactness;
\cite[Lemma~10.4.5 and Definition~10.4.4]{XiaBook} prove independence.
This chapter uses this construction on its positive-ideal-volume locus.

\begin{theorem}\label{thm:ok-transcendental-package}
On the positive multiplier-ideal-volume locus the construction above has
the following
properties.
\begin{enumerate}
\item
\[
 \lambda_n(\Delta_{Y_\bullet}(T))
       =\frac1{n!}\vol(T).
\tag{55.10}\label{eq:ok-trans-volume}
\]
\item If $T\preceq_{\cI}S$, then
$\Delta_{Y_\bullet}(T)\subseteq\Delta_{Y_\bullet}(S)$.
\item On the positive raw-mass locus $\PSH(X,\theta)_{>0}$, the body
map is continuous for $d_S$ and Hausdorff distance.
For arbitrary positive-ideal-volume inputs, convergence of their normalized
$\mathcal I$-envelopes in $d_S$ likewise gives convergence of bodies.
\item If $\pi:Y\to X$ is a modification and $Z_\bullet$ lifts
$Y_\bullet$ with change matrix $g$, then
\[
             \Delta_{Z_\bullet}(\pi^*T)
             =\Delta_{Y_\bullet}(T)g.
\tag{55.11}\label{eq:ok-birational-body}
\]
\item For positive currents $T,S$,
\[
 \Delta_{Y_\bullet}(T)+\Delta_{Y_\bullet}(S)
       \subseteq\Delta_{Y_\bullet}(T+S),
\tag{55.12}\label{eq:ok-minkowski}
\]
with the evident interpretation when their classes differ.
\item For $\lambda>0$,
\[
 \Delta_{Y_\bullet}(\lambda T)=\lambda\Delta_{Y_\bullet}(T).
\tag{55.12a}\label{eq:ok-body-homogeneity}
\]
\end{enumerate}
\par
\end{theorem}

\begin{proof}[Transcendental source boundary]
The analytic-singularity construction, its independence of the lifted
flag, and the K\"ahler-current approximation theorem are imported from
\cite[Definitions~10.4.2--10.4.4, Lemmas~10.4.1--10.4.5 and
Theorem~10.4.2]{XiaBook}.  On a common log resolution they reduce to the ordinary
transcendental Okounkov theorem, translation by divisor valuations, and the
correction matrix for a lifted flag.  The correction-matrix invariance is
the nonformal point that makes \eqref{eq:ok-trans-analytic} independent of
the resolution.

For a general current, the bodies in
\eqref{eq:ok-trans-general} are nested.  Continuity of
$\mathcal I$-model volume under $j^{-1}\omega\downarrow0$, together
with continuity of volume for nested convex bodies, gives
\eqref{eq:ok-trans-volume}.  The ideal-volume perturbation theorem is
\cite[Proposition~7.3.1]{XiaBook}.  Applying the order theorem before taking the
intersection gives item~2.

For continuity on the positive raw-mass locus, ideal projection is
$d_S$-continuous by \cite[Theorem~6.2.3]{XiaBook} and does not change
bodies.  Replace inputs by these normalized representatives.
A $d_S$-convergent sequence can be sandwiched, after passing to a
subsequence, between increasing and decreasing sequences with the same
limit.
The order statement traps the bodies.
Their volumes converge by continuity of multiplier-ideal volume, and equality of
the limiting volume excludes a proper convex subset.

The Minkowski inclusion, including the almost-semigroup estimates needed
already for analytic singularities, and its passage through
quasi-equisingular approximation are part of the imported package from
\cite[Theorem~10.4.2(6)]{XiaBook}.  On the algebraic locus its
counterpart is \cite[Proposition~5.13]{Xia21}.
It is not a formal multiplication argument: multiplier-ideal subadditivity
has the reverse inclusion from the one such an argument would require.
The correction-matrix identity in the imported analytic-singularity
construction proves \eqref{eq:ok-birational-body}; approximation then
gives the stated birational rule.
Finally, current valuations are homogeneous by
\cite[Proposition~10.2.2]{XiaBook}.  In the closed valuative description
of \cite[Theorem~10.4.4]{XiaBook}, scaling the competing currents by
$\lambda$ is a bijection between the competitors for $T$ and $\lambda T$;
ideal order is unchanged after reindexing all positive exponents.
Scaling commutes with closure, proving \eqref{eq:ok-body-homogeneity}.
\end{proof}

\section{Okounkov test curves}
\label{sec:ok-test-curves}

Let $\Delta\subset\mathbf R^n$ be a convex body with nonempty interior.

\begin{definition}[Okounkov test curve]
\label{def:ok-test-curve}
An Okounkov test curve relative to $\Delta$ has a finite terminal
parameter $\Delta_{\max}\in\mathbf R$ and a decreasing family of
nonempty compact convex sets
\[
                  (\Delta_\tau)_{\tau<\Delta_{\max}}
\]
such that
\begin{enumerate}
\item
\[
 t\Delta_\tau+(1-t)\Delta_\sigma
       \subseteq\Delta_{t\tau+(1-t)\sigma}
\quad(0<t<1);
\tag{55.13}\label{eq:ok-body-concavity}
\]
\item $\Delta_\tau\to\Delta$ in Hausdorff distance as
$\tau\to-\infty$.
\end{enumerate}
It is bounded if $\Delta_\tau=\Delta$ for all sufficiently negative
$\tau$.
\end{definition}

The published \cite[Definitions~7.2 and~7.4]{Xia21} additionally requires
finite energy.  Here that is a separate condition, as in
\cite[Definitions~10.5.1--10.5.2]{XiaBook}, so unbounded curves with
infinite negative energy are included.
Concavity implies that $\Delta_\tau$ has positive volume and varies
continuously for $\tau<\Delta_{\max}$.
To prove positivity at $\tau_0$, first choose $\tau_2<\tau_0$ with
positive-volume body, using the initial Hausdorff limit, and then choose
$\tau_1\in(\tau_0,\Delta_{\max})$.
If $\tau_0=(1-t)\tau_2+t\tau_1$, concavity gives volume at least
$(1-t)^n\lambda_n(\Delta_{\tau_2})>0$.
Brunn--Minkowski now makes
$\tau\mapsto\log\lambda_n(\Delta_\tau)$ finite and concave, hence
continuous.  The nested left and right limits are trapped by
$\Delta_\tau$ and have its positive volume; strict convex-volume
monotonicity makes both equal to it.
This is \cite[Proposition~7.3]{Xia21}, or
\cite[Proposition~10.5.1]{XiaBook}.
We retain the closed endpoint
$\Delta_{\Delta_{\max}}:=\bigcap_{\tau<\Delta_{\max}}\Delta_\tau$,
which is nonempty by compactness and can have zero volume.

\begin{definition}[Test function]
\label{def:ok-test-function}
The test function of $\Delta_\bullet$ is
\[
 G_{\Delta_\bullet}(a)
 :=\sup\{\tau<\Delta_{\max}:a\in\Delta_\tau\},
 \qquad a\in\Delta,
\tag{55.14}\label{eq:ok-legendre-function}
\]
with value $-\infty$ if the set is empty.
\end{definition}

\begin{proposition}[Convex test-curve duality]
\label{prop:ok-test-function-duality}
The function $G_{\Delta_\bullet}$ is concave, upper semicontinuous, and
finite on $\Int\Delta$.
Conversely every $G:\Delta\to[-\infty,\infty)$ with those properties
gives an Okounkov test curve by
\[
 \Delta[G]_\tau=\{a\in\Delta:G(a)\geq\tau\},
 \qquad \tau<\sup_\Delta G.
\tag{55.15}\label{eq:ok-inverse-legendre}
\]
The two constructions are inverse, and bounded curves correspond exactly
to test functions bounded from below.
Finite energy is equivalently
\[
                         \int_\Delta G\,\rd\lambda_n>-\infty.
\tag{55.16}\label{eq:ok-finite-energy}
\]
More precisely, their unnormalized energy is
\[
 \begin{split}
 \mathbf E(\Delta_\bullet)
 &:=n!\left(\Delta_{\max}\lambda_n(\Delta)
       +\int_{-\infty}^{\Delta_{\max}}
          (\lambda_n(\Delta_\tau)-\lambda_n(\Delta))\,\rd\tau\right)\\
 &=n!\int_\Delta G\,\rd\lambda_n.
 \end{split}
\tag{55.16a}\label{eq:ok-body-energy}
\]
\end{proposition}

\begin{proof}
The increasing union of the interior-parameter bodies is convex and dense
in $\Delta$, so it contains $\Int\Delta$.  For example, surround any
interior point by a simplex in $\Delta$ and approximate its vertices by
points in that union; a small approximation still surrounds the point.
A common sufficiently negative level contains all the approximating vertices.
Thus $G$ is finite on the interior and bounded above by $\Delta_{\max}$.
Interior-parameter continuity gives
\[
 \{G\geq\tau\}=\Delta_\tau\quad(\tau<\Delta_{\max}),\qquad
 \{G\geq\Delta_{\max}\}=\bigcap_{\tau<\Delta_{\max}}\Delta_\tau.
\tag{55.16b}\label{eq:ok-test-superlevels}
\]
These closed superlevel sets prove upper semicontinuity.  Taking two levels
strictly below the corresponding finite $G$-values and using body concavity
proves concavity of $G$ by passage to the limit.  If either value is
$-\infty$, the concavity inequality is automatic.

Conversely, upper semicontinuity and compactness give a finite maximum of
$G$.  Its nonempty superlevel sets are compact, convex and decreasing;
concavity of $G$ gives Minkowski concavity of the sets.  Their union is dense
because $G$ is finite in the interior.  Hence they have the required initial
Hausdorff limit, and the preceding positivity argument applies.  Formula
\eqref{eq:ok-test-superlevels} proves both inverse identities and the
boundedness assertion.  The finite-energy case is \cite[Theorem~7.5]{Xia21};
the broader test-function interface is \cite[Theorem~10.5.1]{XiaBook}.

Finally, apply Tonelli to the nonnegative identity
\[
 \Delta_{\max}-G(a)
 =\int_{-\infty}^{\Delta_{\max}}
       \mathbf1_{\{G(a)<\tau\}}\,\rd\tau.
\]
It gives \eqref{eq:ok-body-energy} as an extended identity in
$[-\infty,\infty)$, without an undefined difference of infinite terms.
\end{proof}

\begin{lemma}\label{lem:ok-i-model-mass-volume}
If $u\in\PSH(X,\theta)$ is $\mathcal I$-model, then
\[
                         \vol(\theta,u)=M_\theta(u).
\]
In particular, every interior level of a positive-mass
$\mathcal I$-model test curve lies in the domain of
\Cref{def:ok-algebraic-partial-body}.
\end{lemma}

\begin{proof}
The defining equality $P_\theta[u]_{\mathcal I}=u$, inserted into
\eqref{eq:complex-current-volume}, gives
\[
 \vol(\theta,u)
 =\int_X\bigl(\theta+\ddc u\bigr)^n
 =M_\theta(u).
\]
The last assertion follows from positivity of the mass of every interior
level, \Cref{lem:rays-test-positive-levels}, together with
\eqref{eq:complex-mass-volume}.  It is not part of
\Cref{def:rays-test-curve}, whose clause~(3) imposes positive mass at the
endpoint $\Gamma_{-\infty}$ alone; the interior statement is deduced from
that and concavity.
\end{proof}

\begin{theorem}\label{thm:ok-analytic-to-convex}
Let $\Gamma=(\Gamma_\tau)$ be a positive-mass
$\cI$-model test curve in a big class $\alpha$.
Then
\[
 \tau\longmapsto
       \Delta_{Y_\bullet}(\theta,\Gamma_\tau)
\tag{55.17}\label{eq:ok-test-curve-from-gamma}
\]
is an Okounkov test curve relative to
$\Delta_{Y_\bullet}(\theta,\Gamma_{-\infty})$.
Boundedness and finite energy are preserved.
\end{theorem}

\begin{proof}
The family decreases by the $\cI$-order theorem.
Concavity of the analytic test curve, the Minkowski inclusion and
\eqref{eq:ok-body-homogeneity} give \eqref{eq:ok-body-concavity}.
As $\tau\to-\infty$, the level potentials converge in $d_S$ to
$\Gamma_{-\infty}$; hence the bodies converge in Hausdorff distance.
Boundedness is immediate when the analytic levels stabilize.
The diagonal volume formula and \eqref{eq:ok-body-energy} give exactly
\[
 \mathbf E(\Delta_{Y_\bullet}(\theta,\Gamma))
 =\mathbf E^{\Gamma_{-\infty}}(\Gamma)
 =n!\int_{\Delta_{Y_\bullet}(\theta,\Gamma_{-\infty})}G\,\rd\lambda_n,
\tag{55.17a}\label{eq:ok-analytic-convex-energy}
\]
where the middle energy is \eqref{eq:rays-test-energy}.
Thus finite energy is preserved with the stated mass normalization.
The source interface is \cite[Theorem~10.5.2]{XiaBook}.
\end{proof}

\section{The Duistermaat--Heckman measure}
\label{sec:ok-dh}

\begin{definition}
\label{def:ok-dh-body}
The Duistermaat--Heckman measure of an Okounkov test curve is
\[
        \DHm(\Delta_\bullet)
           :=(G_{\Delta_\bullet})_*\!
             \left(\lambda_n|_\Delta\right).
\tag{55.18}\label{eq:ok-dh-body}
\]
It is a Radon measure because $G_{\Delta_\bullet}$ is finite on the
interior and the boundary has Lebesgue measure zero.
\end{definition}

\begin{example}[The interval body of $\mathbf P^1$]
\label{ex:ok-dh-p1-interval}
Take $X=\mathbf P^1$, $L=\mathcal O(1)$, and use the flag given by a
point $p$.  The orders at $p$ of nonzero sections of $L^k$ are
$0,1,\ldots,k$, so the normalized Okounkov body is
\[
                         \Delta=[0,1].
\]
Let $G(x)=x$.  By \eqref{eq:ok-inverse-legendre}, its body curve is
\[
 \Delta_\tau=
 \begin{cases}
 [0,1],&\tau\leq0,\\
 [\tau,1],&0<\tau<1,
 \end{cases}
 \qquad \Delta_{\max}=1.
\]
The function is bounded below, hence the curve is bounded and has finite
energy.  Its Duistermaat--Heckman measure is
\[
                 G_*(\lambda_1|_{[0,1]})
                    =\mathbf 1_{[0,1]}\,\rd x.
\]
In particular, its mass is one and its $m$-th moment is
\[
                 \int_0^1x^m\,\rd x=\frac1{m+1}.
\]
The same number follows directly from \eqref{eq:ok-dh-moment}:
\[
 1+m\int_0^1\tau^{m-1}\bigl((1-\tau)-1\bigr)\,\rd\tau
       =\frac1{m+1}.
\]
This checks the unnormalized Lebesgue-mass convention in dimension one.
\end{example}

\begin{proposition}\label{prop:ok-dh-body-moments}
For every integer $m\geq1$,
\[
\begin{split}
 \int_{\mathbf R}x^m\,\rd \DHm(\Delta_\bullet)(x)
  ={}&\Delta_{\max}^m\lambda_n(\Delta)\\
     &+m\int_{-\infty}^{\Delta_{\max}}
       \tau^{m-1}
       \bigl(\lambda_n(\Delta_\tau)-\lambda_n(\Delta)\bigr)\,\rd \tau,
\end{split}
\tag{55.19}\label{eq:ok-dh-moment}
\]
and
\[
             \DHm(\Delta_\bullet)(\mathbf R)=\lambda_n(\Delta).
\tag{55.20}\label{eq:ok-dh-mass}
\]
Equation \eqref{eq:ok-dh-moment} is an extended-real identity.
\end{proposition}

\begin{proof}
The mass identity is immediate from pushforward.
If $G=G_{\Delta_\bullet}$ is bounded below, Fubini's theorem and
$\{G\geq\tau\}=\Delta_\tau$ give
\[
 G(a)^m=\Delta_{\max}^m
        -\int_{G(a)}^{\Delta_{\max}}m\tau^{m-1}\,\rd \tau.
\]
Integrating yields \eqref{eq:ok-dh-moment}.

For general $G$, put $G_A=\max\{G,-A\}$.
If $m$ is even, $G_A^m\uparrow G^m$ after $A$ passes
$\Delta_{\max}$ in absolute value.
If $m$ is odd, apply monotone convergence to
$\Delta_{\max}^m-G_A^m$.
This proves the claimed extended interpretation without subtracting
undefined quantities of the form $\infty-\infty$.
\end{proof}

\begin{example}[A finite measure with infinite moments]
\label{ex:ok-dh-infinite-energy}
On $\Delta=[0,1]$, set $G(x)=-1/x$ for $x>0$ and $G(0)=-\infty$.
It is concave and upper semicontinuous, with maximum $-1$.
Its curve and measure are
\[
 \Delta_\tau=[-1/\tau,1]\quad(\tau<-1),\qquad
 \DHm(\Delta_\bullet)=\mathbf1_{(-\infty,-1]}(u)\frac{\rd u}{u^2}.
\tag{55.20a}\label{eq:ok-infinite-energy-density}
\]
The density follows from the change of variables $x=-1/u$.
It integrates to one, so it is a finite Radon measure.
For every $m\geq1$, its $m$-th moment is $-\infty$ for odd $m$ and
$+\infty$ for even $m$.
For example, the first-moment formula reads
\[
 -1+\int_{-\infty}^{-1}
          (\lambda_1(\Delta_\tau)-1)\,\rd\tau
 =-1+\int_{-\infty}^{-1}\frac{\rd\tau}{\tau}=-\infty.
\tag{55.20b}\label{eq:ok-infinite-energy-moment}
\]
This also computes its energy as $-\infty$.  Finiteness of the measure
therefore imposes neither finite energy nor finite moments.
\end{example}

\begin{definition}[DH measure of a transcendental metric]
\label{def:ok-dh-metric}
If $X$ admits a smooth flag and $\Gamma$ is a positive-mass
$\mathcal I$-model test curve in the fixed big class, define
\[
 \DHm(\Gamma)
 :=n!\,\DHm\!\left(
       \tau\mapsto\Delta_{Y_\bullet}(\theta,\Gamma_\tau)
                  \right).
\tag{55.21}\label{eq:ok-dh-metric}
\]
\end{definition}

Here a positive-mass test curve is regarded as a transcendental metric via
\eqref{eq:tna-positive-embedding}; consequently
\[
 \vol(\Gamma)
 =M_\theta(\Gamma_{-\infty})
 =\vol(\theta,\Gamma_{-\infty}),
\]
where the last equality is
\Cref{lem:ok-i-model-mass-volume}.

\begin{theorem}\label{thm:ok-dh-flag-independent}
The measure $\DHm(\Gamma)$ is independent of the smooth flag.
It has mass
\[
                         \DHm(\Gamma)(\mathbf R)=\vol(\Gamma),
\tag{55.22}\label{eq:ok-dh-total-mass}
\]
and for $m\geq1$,
\[
\begin{split}
 \int_{\mathbf R}x^m\,\rd \DHm(\Gamma)(x)
 ={}&\Gamma_{\max}^m\vol(\Gamma)\\
 &+m\int_{-\infty}^{\Gamma_{\max}}
 \tau^{m-1}\!
 \left(\int_X\langle\theta_{\Gamma_\tau}^n\rangle
                  -\vol(\Gamma)\right)\,\rd\tau .
\end{split}
\tag{55.23}\label{eq:ok-intrinsic-moments}
\]
\end{theorem}

\begin{proof}
By \Cref{lem:ok-i-model-mass-volume}, the volume formula gives
\[
 n!\lambda_n(\Delta_{Y_\bullet}(\theta,\Gamma_\tau))
 =M_\theta(\Gamma_\tau)\qquad(\tau<\Gamma_{\max}).
\]
Substitution in \eqref{eq:ok-dh-moment} proves the moment identity, and
its initial-level version proves the total mass.
More directly, \eqref{eq:ok-test-superlevels} identifies every tail:
\[
 \DHm(\Gamma)([\tau,\infty))=M_\theta(\Gamma_\tau)
       \quad(\tau<\Gamma_{\max}),\qquad
 \DHm(\Gamma)((\Gamma_{\max},\infty))=0.
\tag{55.23a}\label{eq:ok-intrinsic-tail}
\]
All these quantities are independent of the flag.  Increasing $\tau$ to
$\Gamma_{\max}$ also determines its possible endpoint atom by continuity
from above of a finite measure.  Half-lines determine a finite Borel measure,
so the tail identity proves flag independence for every positive-mass curve,
including curves with infinite moments or without finite energy.
The corresponding source package is
\cite[Theorem~13.3.1]{XiaBook}; the tail argument here gives the
measure-uniqueness step directly.
\end{proof}

If a compact K\"ahler manifold has no smooth flag, choose a modification
$\pi:Y\to X$ admitting one and set
\[
                         \DHm(\Gamma):=\DHm(\pi^*\Gamma).
\tag{55.24}\label{eq:ok-dh-modification}
\]
The birational rule \eqref{eq:ok-birational-body} and flag independence make
this independent of $\pi$.
Thus \eqref{eq:ok-dh-modification} also proves birational invariance.

\begin{proposition}\label{prop:ok-dh-weak-continuity}
Let $\Gamma^j,\Gamma$ be positive-mass $\mathcal I$-model test curves
in one fixed big background $\theta$.
Then
\[
                    \DHm(\Gamma^j)\rightharpoonup\DHm(\Gamma)
\tag{55.25}\label{eq:ok-dh-weak}
\]
in either of the following cases:
\begin{enumerate}
\item $\Gamma^j\searrow\Gamma$ and
$\vol(\Gamma^j)\to\vol(\Gamma)$;
\item $\Gamma^j\nearrow\Gamma$, where the limit is the upper
semicontinuous supremum in the test-curve space.
\end{enumerate}
\par
\end{proposition}

\begin{proof}
All curves lie in one fixed big background $\theta$ and have normalized
ideal-model levels.  In the decreasing case the terminal parameters decrease
to $\Gamma_{\max}$ and the interior levels decrease to $\Gamma_\tau$.
The positive mass of this limit permits
\cite[Proposition~3.1.12 and Corollary~6.2.6]{XiaBook}, giving convergence
of each interior-level mass.  The hypothesis on total volumes is also needed:
initial levels need not commute with the infimum.
Indeed put $\eta=\inf_j\Gamma^j_{-\infty}$.  It dominates
$\Gamma_{-\infty}$ and is a normalized model potential whose mass is
$\lim_j\vol(\Gamma^j)>0$, by the same decreasing-model theorem.
Equality of this mass with $M_\theta(\Gamma_{-\infty})$ and domination
for model potentials, \cite[Proposition~6.1.1]{XiaBook}, give
$\eta=\Gamma_{-\infty}$.

In the increasing case, the terminal parameters increase to
$\Gamma_{\max}$.  At each fixed interior parameter, the levels are eventually
nontrivial and increase after upper semicontinuous regularization to
$\Gamma_\tau$.  Initial levels have the same increasing-supremum rule.
Thus increasing mass continuity, \Cref{prop:complex-mass-continuity}, gives
both the level-mass limits and $\vol(\Gamma^j)\to\vol(\Gamma)$.
These are the fixed-class monotone-limit interfaces of
\Cref{prop:rays-test-limits}; equivalently see
\cite[Proposition~13.3.1]{XiaBook}.

The intrinsic tail formula \eqref{eq:ok-intrinsic-tail} now gives convergence
of the distributions at every interior parameter.  Their tails are zero
above the terminal parameters, which converge as just stated.
Interior-level mass is continuous in the parameter by body continuity, so
there are no interior atoms; the terminal parameter may be an atom and is
omitted from the distribution-function convergence criterion.
Together with convergence of the total masses, that criterion yields weak
convergence of finite measures.  To see explicitly that mass does not escape,
all right endpoints are bounded by one constant.  At a sufficiently negative
$\tau$, the limit's missing tail mass
$\vol(\Gamma)-M_\theta(\Gamma_\tau)$ is arbitrarily small; the same is true
for the sequence by the level and total-mass limits.
This gives tightness and hence convergence against all bounded continuous
test functions, not only compactly supported ones.
\end{proof}

For an algebraic test configuration of an ample line bundle $L$, with
weights $w_{m,j}$ on $H^0(X,L^m)$ in the parameter convention of
\Cref{ex:rf-p1-complete}, the intersection-mass convention here is
\[
 \DHm(\Gamma)
 =\lim_{m\to\infty}\frac{n!}{m^n}\sum_j\delta_{w_{m,j}/m}.
\tag{55.25a}\label{eq:ok-dh-weight-normalization}
\]
The probability distribution is $\DHm(\Gamma)/(L^n)$.
The published \cite[Definitions~7.6--7.8, Theorem~7.9 and
Remark~7.17]{Xia21} uses the body-level measure, of mass $(L^n)/n!$.
Its comparison with \cite[Section~3.2]{BHJ17} must therefore be read after
this mass conversion.

\begin{example}[The product square]
\label{ex:ok-dh-product-square}
On $X=\mathbf P^1\times\mathbf P^1$, take $L=\mathcal O(1,1)$ and the
product of the standard actions of \Cref{ex:rf-p1-complete}.
At degree $m$, monomials indexed by $(i,j)\in\{0,\ldots,m\}^2$ have
weights $i+j$.  The toric flag gives $\Delta=[0,1]^2$ and $G(x,y)=x+y$:
the filtered normalized lattice sets converge to the square's superlevel
sets.  Their areas are
\[
 \lambda_2(\Delta_\tau)=
 \begin{cases}
 1,&\tau\leq0,\\
 1-\tau^2/2,&0\leq\tau\leq1,\\
 (2-\tau)^2/2,&1\leq\tau\leq2.
 \end{cases}
\tag{55.25b}\label{eq:ok-dh-square-profile}
\]
Thus the metric measure, including the factor $2!$, is
\[
 \DHm(\Gamma)=
 \begin{cases}
 2x\,\rd x,&0\leq x\leq1,\\
 2(2-x)\,\rd x,&1\leq x\leq2,\\
 0,&\text{elsewhere}.
 \end{cases}
\tag{55.25c}\label{eq:ok-dh-square-density}
\]
Its mass is $(L^2)=2$, its first moment is $2$, and its second moment is
$7/3$.  Division by two gives the probability law of a sum of two independent
uniform variables.  Formula \eqref{eq:ok-dh-weight-normalization} gives the
same law by two-dimensional lattice Riemann sums.
\end{example}

\section{Mixed volumes of partial Okounkov bodies}
\label{sec:ok-mixed-volumes}
\index{Okounkov body!mixed volume}\index{movable product}

Everything so far attached one convex body to one current and matched its
Euclidean volume with the volume of its $\mathcal I$-model envelope.  We
now prove the corresponding mixed formula.  Throughout this section $X$
is a connected smooth projective variety over $\mathbf C$ of dimension
$n$, and $\nu$ ranges over the surjective valuations
\[
 \nu:\mathbf C(X)^\times\longrightarrow\mathbf Z^n_{\mathrm{lex}}
 \quad\text{trivial on }\mathbf C^\times ,
\tag{55.26}\label{eq:ok-mvob-valuations}
\]
where $\mathbf Z^n_{\mathrm{lex}}$ carries the lexicographic order.  Flag
valuations form a subfamily, but every supremum in this section is over the
whole family in \eqref{eq:ok-mvob-valuations}.

\begin{lemma}\label{lem:ok-mvob-valuative-leaves}
Let $\nu$ be as in \eqref{eq:ok-mvob-valuations}.  Its residue field is
$\mathbf C$, every nonzero leaf of its valuation filtration is
one-dimensional, and for every nonzero finite-dimensional
$W\subset\mathbf C(X)$,
\[
                         \#\nu(W^\times)=\dim_{\mathbf C}W.
\tag{55.27}\label{eq:ok-mvob-leaf-count}
\]
The valuation has a center on every proper model of $\mathbf C(X)$.
Changing a rational trivialization of a line bundle translates its
normalized partial body and hence leaves all mixed volumes unchanged.
\end{lemma}

\begin{proof}
The rational rank of $\nu$ is $n$.  Abhyankar's inequality therefore
forces the residue transcendence degree to be zero.  The residue field is
an algebraic extension of the algebraically closed field $\mathbf C$, so
it is $\mathbf C$.  After division by any element of value $a$, the
nonzero leaf
$\mathcal F_{\succeq a}/\mathcal F_{\succ a}$ embeds in the residue
field.  It is consequently one-dimensional.
Vectors with distinct values are linearly independent, since a finite
nonzero linear combination has a unique least-valued summand.
Consequently there are at most $\dim W$ values, so the value set is finite.
Filtering $W$ by these values and summing its one-dimensional leaves
proves \eqref{eq:ok-mvob-leaf-count}.  The existence of a center
is the valuative criterion of properness.

If rational frames $e,e'$ satisfy $e'=ge$, the functions representing
sections of $L^{\otimes k}$ differ by multiplication by $g^{-k}$.
Their normalized valuation sets therefore differ by the fixed translation
$-\nu(g)$.  Translation invariance of mixed volume proves the last
assertion.
\end{proof}

\begin{remark}[Two normalizations that must be converted, not assumed]
\label{rem:ok-mvob-normalizations}
Two conversions are needed before the source statements can be read against
the conventions of this book.

\emph{Mixed volume.}  For convex bodies $K_1,\dots,K_n\subset\mathbf R^n$
this book writes $V(K_1,\dots,K_n)$ for the mixed volume normalized by
$V(K,\dots,K)=\lambda_n(K)$, so that $n!\,V(K_1,\dots,K_n)$ --- not
$V(K_1,\dots,K_n)$ --- is the coefficient of $t_1\cdots t_n$ in
$\lambda_n(\sum_it_iK_i)$, as in
\Cref{thm:appD-minkowski-polynomial}.  (Setting every $K_i=K$ checks the
factor: the coefficient of $t_1\cdots t_n$ in
$(\sum_it_i)^n\lambda_n(K)$ is $n!\,\lambda_n(K)$.)  The source writes $\vol$ for the
same quantity.  We use $V$.

\emph{Volume of a metric.}  The source's volume of a psh metric is the one
detected by multiplier ideals, that is
$\vol(\theta,\varphi)=\int_X(\theta+\ddc P_\theta[\varphi]_{\cI})^n$ of
\Cref{def:complex-current-volume}, and not the bare non-pluripolar mass
$M_\theta(\varphi)$ of \eqref{eq:complex-npp-mass}.  By
\eqref{eq:complex-mass-volume} one has
$M_\theta(\varphi)\leq\vol(\theta,\varphi)$, with equality for analytic
singularities and, more generally, for $\mathcal I$-model potentials.
The hypotheses of \Cref{def:ok-algebraic-partial-body} and
\Cref{thm:ok-transcendental-package} are therefore exactly the
positive-volume hypothesis appearing in \eqref{eq:ok-mvob-metrics}.  The
mixed quantity $\vol(\ddc h_1,\dots,\ddc h_n)$ is the mixed
non-pluripolar mass of the $\cI$-model envelopes,
\[
 \vol(\ddc h_1,\dots,\ddc h_n)
 =\int_X\bigl(\theta_1+\ddc P_{\theta_1}[\varphi_1]_{\cI}\bigr)
        \wedge\cdots\wedge
        \bigl(\theta_n+\ddc P_{\theta_n}[\varphi_n]_{\cI}\bigr),
\tag{55.28}\label{eq:ok-mvob-mixed-mass}
\]
where $h_i=h_{0,i}\rme^{-\varphi_i}$ as in
\eqref{eq:ok-partial-metric} and the wedge is the non-pluripolar product
\eqref{eq:complex-global-npp}.
\end{remark}

\subsection{Order and the algebraic upper bound}
\label{sec:ok-mixed-upper}

\begin{proposition}
\label{prop:ok-mvob-superadditivity}
Let $L,L'$ be big line bundles on $X$, let $h,h'$ be psh metrics on
$L,L'$, and let $\nu$ satisfy \eqref{eq:ok-mvob-valuations}.  Assume
that all the partial bodies below are defined, that is, that the relevant
potentials have positive multiplier-ideal volume.  Then
\[
 \Delta_\nu(L,h)+\Delta_\nu(L',h')
   \subseteq\Delta_\nu(L\otimes L',h\otimes h'),
\tag{55.29}\label{eq:ok-mvob-superadditivity}
\]
and, for one fixed smooth background $\theta$ and
$\varphi,\psi\in\PSH(X,\theta)$,
\[
 \varphi\preceq_{\cI}\psi
   \quad\Longrightarrow\quad
 \Delta_\nu(\theta,\varphi)\subseteq\Delta_\nu(\theta,\psi).
\tag{55.30}\label{eq:ok-mvob-i-monotone}
\]
\end{proposition}

\begin{proof}
The general-valuation Minkowski inclusion is
\cite[Proposition~10.3.10]{XiaBook}; it gives
\eqref{eq:ok-mvob-superadditivity}.  A naive levelwise product argument
would not suffice, since subadditivity gives
$\cI(k(\varphi+\varphi'))
 \subseteq\cI(k\varphi)\cI(k\varphi')$, the opposite inclusion.
The implication \eqref{eq:ok-mvob-i-monotone} follows directly from
\eqref{eq:ok-i-order} of \Cref{thm:ok-algebraic-volume-order}.
\end{proof}

\begin{lemma}\label{lem:ok-mvob-finite-level}
Let $F/\mathbf C$ be finitely generated of transcendence degree $n$,
let $X_0$ be an integral projective model of $F$, and let $\nu$ be as
in \eqref{eq:ok-mvob-valuations}.  Let
$0\ne W_i\subset F$ be finite-dimensional, let $Y$ be the closure of
the image of the joint evaluation map
\[
 X_0\dashrightarrow\prod_{i=1}^n\mathbf P(W_i^\vee),
\]
and suppose that this map is birational onto $Y$.  If
\[
 \operatorname{span}_{\mathbf Z}
 \{\nu(f)-\nu(g):f,g\in W_i^\times,\ 1\leq i\leq n\}
 =\mathbf Z^n,
\tag{55.31}\label{eq:ok-mvob-finite-level-degree}
\]
then, with $P_i=\operatorname{conv}\nu(W_i^\times)$ and with $H_i$
the hyperplane class from the $i$-th factor,
\[
 n!V(P_1,\ldots,P_n)\leq(H_1\cdots H_n\cdot Y).
\tag{55.32}\label{eq:ok-mvob-admissible}
\]
\end{lemma}

\begin{proof}
By \Cref{lem:ok-mvob-valuative-leaves}, choose a basis $f_{ij}$ of each
$W_i$, one vector for each value $a_{ij}=\nu(f_{ij})$.  Choose
$u_1,\ldots,u_n\in F^\times$ with $\nu(u_\ell)=e_\ell$.  The residue
field calculation in that lemma gives a graded isomorphism
\[
 \mathbf C[\mathbf Z^n]\simeq\operatorname{gr}_\nu F,
 \qquad \chi^a\longmapsto\operatorname{in}_\nu(u^a).
\]
After rescaling the $f_{ij}$, their initial forms are $\chi^{a_{ij}}$.

Let $S=\mathbf C[x_{ij}]$, multigraded by $\deg x_{ij}=e_i$, and let
$I$ be the multihomogeneous prime ideal of $Y$.  Give $x_{ij}$ the
lexicographic matrix weight $a_{ij}$, set
$Q=\operatorname{in}_{\mathsf A}(I)$, and put
\[
 J=\ker\!\left(S\longrightarrow
 \mathbf C[\mathbf Z^n][t_1,\ldots,t_n],\quad
 x_{ij}\longmapsto\chi^{a_{ij}}t_i\right).
\]
Taking the least-valuation part of each relation among the $f_{ij}$ gives
only the inclusion
\[
                              Q\subseteq J.
\tag{55.33}\label{eq:ok-mvob-initial-inclusion}
\]
By \cite[Proposition~8.10]{KM19}, a rational scalar weight realizes the
matrix initial ideal; clear its denominators.  The resulting flat
multigraded degeneration preserves every multihomogeneous piece:
\[
 \dim(S/I)_{\mathbf d}=\dim(S/Q)_{\mathbf d}
 \qquad(\mathbf d\in\mathbf N^n).
\tag{55.34}\label{eq:ok-mvob-hilbert-functions}
\]

The ideal $J$ defines the closure $Z$ of the monomial map
\[
 (\mathbf C^\times)^n\longrightarrow
 \prod_i\mathbf P(W_i^\vee),\qquad
 z\longmapsto([z^{a_{i0}}:\cdots:z^{a_{ir_i}}])_i.
\]
Condition \eqref{eq:ok-mvob-finite-level-degree} says that the exponent
differences generate the full character lattice.  Thus this map is a
closed immersion on the torus and $\dim Z=n$.  By
\eqref{eq:ok-mvob-initial-inclusion}, $Z$ lies in
$Y_0=\operatorname{MultiProj}(S/Q)$.  Passing to
$\operatorname{MultiProj}$ removes only multigraded irrelevant torsion;
such torsion vanishes in all sufficiently large multidegrees, and hence
does not change the eventual Hilbert polynomial.  Equation
\eqref{eq:ok-mvob-hilbert-functions} therefore gives
$\dim Y_0=\dim Y=n$, so $Z$ is a top-dimensional component of $Y_0$.
Writing the top-dimensional cycle as
$[Y_0]_n=\sum_Wm_W[W]$, positivity of the multiplicities and nefness of
the globally generated $H_i$ give
\[
 (H_1\cdots H_n\cdot Y)
 =(H_1\cdots H_n\cdot[Y_0]_n)
 \geq(H_1\cdots H_n\cdot Z).
\tag{55.35}\label{eq:ok-mvob-toric-component}
\]

If some $W_i$ is one-dimensional, then $P_i$ is a point, so the mixed
volume in \eqref{eq:ok-mvob-admissible} is zero and the result follows.
Otherwise, for a fixed point of
$Z\setminus(\mathbf C^\times)^n$, requiring one hyperplane in each factor
to pass through it has total codimension $n$.  Since the boundary has
dimension at most $n-1$, the incidence correspondence shows that general
hyperplanes have no common boundary point.  Their intersection is computed
on the dense torus.  The exponent differences generate $\mathbf Z^n$, so
the Bernstein--Kushnirenko formula \cite{Ber75} has no lattice-index factor
and yields
\[
                    (H_1\cdots H_n\cdot Z)=n!V(P_1,\ldots,P_n).
\]
Together with \eqref{eq:ok-mvob-toric-component}, this proves
\eqref{eq:ok-mvob-admissible}.
\end{proof}

\begin{proposition}\label{prop:ok-mvob-movable-upper}
For big $\mathbf Q$-divisors $D_1,\ldots,D_n$ and every valuation in
\eqref{eq:ok-mvob-valuations},
\[
 n!V(\Delta_\nu(D_1),\ldots,\Delta_\nu(D_n))
 \leq\langle D_1,\ldots,D_n\rangle.
\tag{55.36}\label{eq:ok-mvob-movable-upper}
\]
\end{proposition}

\begin{proof}
By homogeneity, clear denominators.  For
$W_{i,k}=H^0(X,\mathcal O_X(kD_i))$, use rational trivializations at the
center of $\nu$; by \Cref{lem:ok-mvob-valuative-leaves}, this only
translates normalized bodies.  For a sufficiently divisible $p$, every
$|pD_i|$ is birational.  To obtain the lattice assertion at one common
degree, use the degree-zero fraction field of the section ring of $D_1$,
which is $\mathbf C(X)$.  Choose rational functions $g_j$ with
$\nu(g_j)=e_j$, and write $g_j=a_j/b_j$ with $a_j,b_j$ homogeneous
sections of the same degree $q$.  Put $s_0=\prod_jb_j$ and
$s_j=a_j\prod_{\ell\ne j}b_\ell$, and set $p_0=nq$: each of these is a
product of $n$ sections of degree $q$, so all $s_j$ lie in one
$W_{1,p_0}$, and $\nu(s_j)-\nu(s_0)=\nu(a_j)-\nu(b_j)=e_j$.  Replace the previously chosen
birational degree by a common multiple $p=rp_0$, and replace these sections by
$s_0^r,s_js_0^{r-1}$.  For $k=mp$, the sections
$s_0^{rm},s_js_0^{rm-1}$ retain the unit differences, while the
symmetric-power subsystem induces the Veronese of each degree-$p$ map.
Thus \Cref{lem:ok-mvob-finite-level} applies at every such $k$.

Resolve all degree-$k$ base ideals on one model and write
\[
 \mu_k^*(kD_i)=M_{i,k}+F_{i,k},\qquad
 M_{i,k}\ \text{basepoint free},\quad F_{i,k}\geq0.
\]
The projection formula, in the form of \Cref{thm:appF-projection}, gives
\[
 n!V(\operatorname{conv}\nu(W_{1,k}^\times),\ldots,
      \operatorname{conv}\nu(W_{n,k}^\times))
 \leq M_{1,k}\cdots M_{n,k}.
\]
The nef approximants $k^{-1}M_{i,k}$ lie below the pullbacks of $D_i$.
The nef-approximation inequality
\cite[Propositions~2.12 and~2.13]{BFJ09} therefore gives
\[
 M_{1,k}\cdots M_{n,k}
 \leq k^n\langle D_1,\ldots,D_n\rangle.
\]
Divide by $k^n$.  The normalized finite-level convex hulls converge in
Hausdorff distance by \cite[Proposition~10.3.4]{XiaBook}; continuity of
mixed volume in \Cref{thm:appD-minkowski-polynomial} proves
\eqref{eq:ok-mvob-movable-upper}.
\end{proof}

\subsection{Uniformity in the valuation}
\label{sec:ok-mixed-uniform}

\begin{lemma}\label{lem:ok-mvob-potential-sandwich}
Let $u_j,u\in\PSH(X,\theta)_{>0}$ be normalized model potentials with
$u_j\to u$ in $d_S$.  After passing to a subsequence there are
increasing normalized model potentials $\eta_j$ and decreasing normalized
model potentials $\psi_j$ such that
\[
 \eta_j\leq u_j\leq\psi_j,\qquad
 \eta_j\leq u\leq\psi_j,\qquad
 \eta_j,\psi_j\longrightarrow u\quad\text{in }d_S.
\tag{55.37}\label{eq:ok-mvob-potential-sandwich}
\]
For finitely many sequences, the same subsequence may be chosen for all of
them.
\end{lemma}

\begin{proof}
The sandwich theorem \cite[Proposition~6.2.3]{XiaBook} supplies the first
pair of inequalities and the two convergences.  If its upper sequence is
not model, replace it by $P_\theta[\psi_j]$.  Monotonicity preserves
decrease and the upper inequality, while model projection preserves the
$d_S$-limit.

Fix $j$.  For $k\geq j$,
$\eta_j\leq\eta_k\leq u_k$.  Closure of the $P$-preorder under
simultaneous $d_S$-limits \cite[Proposition~6.2.5]{XiaBook}, applied to
the constant sequence $\eta_j$ and the tail $(u_k)_{k\geq j}$, gives
$\eta_j\preceq_Pu$.  Since both potentials are normalized and model,
this is $\eta_j\leq u$.  Similarly,
$u_k\leq\psi_k\leq\psi_j$ for $k\geq j$, so order closure and modelness
give $u\leq\psi_j$.  This proves the missing two inequalities.  For
finitely many sequences, apply the construction successively and take the
finite diagonal subsequence.
\end{proof}

\par\noindent\begin{minipage}{\linewidth}
\begin{lemma}\label{lem:ok-mvob-uniform-sandwich}
For $1\leq i\leq n$, suppose that nonempty compact convex sets satisfy
all four inclusions
\[
 A_{i,j}\subseteq C_{i,j}\subseteq B_{i,j},\qquad
 A_{i,j}\subseteq D_{i,j}\subseteq B_{i,j}.
\tag{55.38}\label{eq:ok-mvob-convex-sandwich}
\]
Assume, for constants $v,M>0$ independent of $i,j$, that
\[
 \lambda_n(A_{i,j})\geq v,\qquad
 \lambda_n(B_{1,j}+\cdots+B_{n,j})\leq M,\qquad
 \lambda_n(B_{i,j})-\lambda_n(A_{i,j})\longrightarrow0.
\]
Then
\[
 \left|V(C_{1,j},\ldots,C_{n,j})
       -V(D_{1,j},\ldots,D_{n,j})\right|\longrightarrow0.
\]
\end{lemma}
\end{minipage}\par


\begin{proof}
For each $i,j$, translate the whole quadruple so that
$0\in A_{i,j}$.  Every body is then contained in
$S_j=B_{1,j}+\cdots+B_{n,j}$, and
$v\leq\lambda_n(S_j)\leq M$.  Send the John ellipsoid of $S_j$ to the
unit ball by an affine map with linear part $L_j$.  John's theorem
\cite[Theorem~10.12.2]{Sch14} gives
\[
 \mathbf B^n\subseteq L_jS_j+b_j\subseteq n\mathbf B^n,\qquad
 \frac{\lambda_n(\mathbf B^n)}{M}\leq|\det L_j|
 \leq\frac{n^n\lambda_n(\mathbf B^n)}{v}.
\]
All transformed bodies lie in one compact set, their inner volumes stay
positive, and their inner--outer volume gaps tend to zero.  From any
subsequence, Blaschke selection \cite[Theorem~1.8.7]{Sch14} gives a further
subsequence on which all
$4n$ bodies converge.  At the limit, the four inclusions, equality of the
positive inner and outer volumes, and strict volume monotonicity under a
proper inclusion of full-dimensional convex bodies imply
$\widehat A_i=\widehat B_i=\widehat C_i=\widehat D_i$.  Continuity of
mixed volume gives the conclusion after transformation.  The lower bound
for $|\det L_j|$ transfers it back.  Since every subsequence has such a
further subsequence, the full sequence converges.
\end{proof}

\begin{proposition}\label{prop:ok-mvob-uniform-ds}
For $1\leq i\leq n$, let $L_i$ be a big line bundle, choose a smooth
form $\theta_i\in c_1(L_i)$, let
$\varphi_{i,j},\varphi_i\in\PSH(X,\theta_i)_{>0}$, and assume that
$\varphi_{i,j}\to\varphi_i$ in $d_S$.
Then
\[
 \sup_\nu\left|
 V(\Delta_\nu(\theta_1,\varphi_{1,j}),\ldots,
   \Delta_\nu(\theta_n,\varphi_{n,j}))
 -V(\Delta_\nu(\theta_1,\varphi_1),\ldots,
   \Delta_\nu(\theta_n,\varphi_n))\right|\longrightarrow0.
\tag{55.39}\label{eq:ok-mvob-uniform-ds}
\]
\end{proposition}

\begin{proof}
Replace all potentials by normalized model envelopes.  If
\eqref{eq:ok-mvob-uniform-ds} failed, an approximate maximizer of the
supremum would give, after a subsequence, valuations $\nu_j$ and
$\delta>0$ for which the displayed difference is at least $\delta$.
Use the common subsequence in \Cref{lem:ok-mvob-potential-sandwich} and set
\[
 \begin{aligned}
 A_{i,j}&=\Delta_{\nu_j}(\theta_i,\eta_{i,j}),&
 B_{i,j}&=\Delta_{\nu_j}(\theta_i,\psi_{i,j}),\\
 C_{i,j}&=\Delta_{\nu_j}(\theta_i,\varphi_{i,j}),&
 D_{i,j}&=\Delta_{\nu_j}(\theta_i,\varphi_i).
 \end{aligned}
\]
Discard a finite initial segment.  Continuity of the diagonal masses and
positivity of the limits then put every one of these bodies on the
positive-volume locus.  The two sandwiches in
\eqref{eq:ok-mvob-potential-sandwich} and
\eqref{eq:ok-mvob-i-monotone} give all four inclusions in
\eqref{eq:ok-mvob-convex-sandwich}.  The diagonal volume formula and
$d_S$-continuity give
\[
 \lambda_n(B_{i,j})-\lambda_n(A_{i,j})\longrightarrow0,\qquad
 \lambda_n(A_{i,j})\longrightarrow
 \frac1{n!}\vol(\theta_i,\varphi_i)>0,
\]
independently of $\nu_j$.  Moreover,
\[
 B_{1,j}+\cdots+B_{n,j}
 \subseteq\Delta_{\nu_j}\!\left(\sum_i\theta_i,
                                  \sum_i\psi_{i,j}\right)
 \subseteq\Delta_{\nu_j}\!\left(\bigotimes_iL_i\right)
\]
by \eqref{eq:ok-mvob-superadditivity} and containment in the complete
series.  The last body has Euclidean volume
$\vol(\bigotimes_iL_i)/n!$, independent of $j$ and $\nu_j$.
\Cref{lem:ok-mvob-uniform-sandwich} now contradicts $\delta$.  This
proves \eqref{eq:ok-mvob-uniform-ds}.  Notice that the conclusion concerns
mixed-volume values uniformly in $\nu$, not Hausdorff convergence
uniformly in $\nu$.
\end{proof}

\subsection{Rational analytic singularities}
\label{sec:ok-mixed-analytic}

\begin{definition}\label{def:ok-rational-analytic}
A psh metric has \emph{rational analytic singularities} if locally its
weight is
\[
 c\log\!\left(\sum_j|f_j|^2\right)+O(1),\qquad c\in\mathbf Q_{>0}.
\]
This is the rational subcase of
\Cref{def:complex-analytic-singularities}; the rationality is needed for
the divisor and section arguments below.
\end{definition}

\begin{lemma}\label{lem:ok-mvob-effective-shift}
Let $A,B$ be big $\mathbf Q$-divisors, let $E\geq0$ be a
$\mathbf Q$-divisor, and suppose
$A\sim_{\mathbf Q}aB+bE$ with $a,b\in\mathbf Q_{\geq0}$.  Then
\[
 a\Delta_\nu(B)+b\nu(E)\subseteq\Delta_\nu(A).
\tag{55.40}\label{eq:ok-mvob-effective-shift}
\]
\end{lemma}

\begin{proof}
Clear the denominators in the divisor relation and in $a,b$.  Multiplying
a section of the corresponding multiple of $B$ by the required power of
the canonical section of $E$ produces a section of the corresponding
multiple of $A$.  Valuation vectors add under multiplication.  Divide by
the common multiple and take the closed convex hull to obtain
\eqref{eq:ok-mvob-effective-shift}.
\end{proof}

\begin{proposition}\label{prop:ok-mvob-analytic-case}
Assume that the psh metrics $h_i$ on the big line bundles $L_i$ have
rational analytic singularities and that their curvature currents are
K\"ahler currents.  Then \eqref{eq:ok-mvob-metrics} holds.
\end{proposition}

\begin{proof}
On a common projective log resolution $\pi:Y\to X$, write
\[
 \pi^*c_1(L_i,h_i)=[D_i]+R_i,\qquad
 A_i:=\pi^*L_i-D_i\in\operatorname{Pic}(Y)_{\mathbf Q},
\]
where $D_i\geq0$ is rational and $R_i$ has bounded potentials.  The
class $\alpha_i=c_1(A_i)=\{R_i\}$ is nef and big.  Birational invariance
of the non-pluripolar product gives
\[
 \int_X\bigwedge_i c_1(L_i,h_i)=\alpha_1\cdots\alpha_n.
\tag{55.41}\label{eq:ok-mvob-resolution-mass}
\]
For every valuation of the common function field, the log-resolution
formula \cite[Remark~10.3.1, equation~(10.22)]{XiaBook} gives
\[
 \Delta_\nu(L_i,h_i)=\Delta_\nu(A_i)+\nu(D_i).
\tag{55.42}\label{eq:ok-mvob-resolution-body}
\]
Translations do not affect mixed volume.  Applying
\Cref{prop:ok-mvob-movable-upper}, and using nefness, yields
\[
 n!V(\Delta_\nu(L_1,h_1),\ldots,\Delta_\nu(L_n,h_n))
 \leq\alpha_1\cdots\alpha_n.
\tag{55.43}\label{eq:ok-mvob-analytic-upper}
\]

For the reverse inequality, fix an ample $\mathbf Q$-divisor $H$ on
$Y$.  Choose one rational $c>0$ such that every $cA_i-H$ is big, and
choose $E_i\geq0$ with $E_i\sim_{\mathbf Q}cA_i-H$.  For rational
$\varepsilon>0$, set $B_i=A_i+\varepsilon H$.  Then $B_i$ is ample
and
\[
 A_i\sim_{\mathbf Q}
 \frac1{1+c\varepsilon}B_i
 +\frac{\varepsilon}{1+c\varepsilon}E_i.
\tag{55.44}\label{eq:ok-mvob-kodaira}
\]
By \cite[Theorem~1.1]{Wil25}, there is an admissible flag, allowed to depend
on $\varepsilon$, whose valuation $\nu_\varepsilon$ satisfies
\[
 n!V(\Delta_{\nu_\varepsilon}(B_1),\ldots,
      \Delta_{\nu_\varepsilon}(B_n))=B_1\cdots B_n.
\tag{55.45}\label{eq:ok-mvob-wilms}
\]
Apply \Cref{lem:ok-mvob-effective-shift} to
\eqref{eq:ok-mvob-kodaira}.  Translation invariance and
\eqref{eq:ok-mvob-wilms} give
\[
 \sup_\nu V(\Delta_\nu(A_1),\ldots,\Delta_\nu(A_n))
 \geq\frac{B_1\cdots B_n}{n!(1+c\varepsilon)^n}.
\tag{55.46}\label{eq:ok-mvob-analytic-lower}
\]
Let rational $\varepsilon\downarrow0$.  Equations
\eqref{eq:ok-mvob-resolution-mass}--\eqref{eq:ok-mvob-analytic-lower}
prove the proposition.  No maximizing valuation was used: at each
$\varepsilon$ a possibly different flag supplies the lower bound.
\end{proof}

\subsection{The positive-volume formula}
\label{sec:ok-mixed-formula}

\begin{lemma}\label{lem:ok-mvob-i-model-approximation}
Let $h=h_0\rme^{-\varphi}$ be a psh metric of positive
multiplier-ideal volume and put $u=P_\theta[\varphi]_{\cI}$.  Then
$u\sim_{\cI}\varphi$, the partial bodies of $u$ and $\varphi$ agree,
$u$ is a model potential with
$M_\theta(u)=\vol(\theta,\varphi)>0$, and there are rational-analytic
potentials $u_j$ such that
\[
 u_j\longrightarrow u\quad\text{in }d_S,\qquad
 \theta+\ddc u_j\ \text{is a K\"ahler current}.
\tag{55.47}\label{eq:ok-mvob-approximating}
\]
\end{lemma}

\begin{proof}
The first two assertions follow from
\Cref{thm:complex-i-projection} and the $\cI$-invariance of partial
bodies.  It is also model: the model envelope
$P_\theta[u]$ is $\cI$-equivalent to $u$, so maximality of the
$\cI$-envelope gives $P_\theta[u]\leq u$, while normalized $u$ is a
competitor for its model envelope and gives the reverse inequality.  The
mass--volume identity for an $\cI$-model potential is
\eqref{eq:complex-mass-volume}.  Thus the assumed ideal volume, rather than
the possibly smaller raw mass of $\varphi$, is the positive mass of
$u$.  The potential $u$ is $\cI$-good, and the rational-analytic
K\"ahler-current approximation in \cite[Theorem~7.1.1]{XiaBook} gives
\eqref{eq:ok-mvob-approximating}.
\end{proof}

\par\noindent\begin{minipage}{\linewidth}
\begin{theorem}\label{thm:ok-mvob-metrics}
Let $L_i$ be big line bundles on $X$, and let $h_i$ be psh metrics
with $\vol(\theta_i,\varphi_i)>0$, in the normalization of
\Cref{rem:ok-mvob-normalizations}.  Then
\[
 \vol(\ddc h_1,\dots,\ddc h_n)
 =n!\,\sup_\nu V\!\left(
    \Delta_\nu(L_1,h_1),\dots,\Delta_\nu(L_n,h_n)\right).
\tag{55.48}\label{eq:ok-mvob-metrics}
\]
\end{theorem}
\end{minipage}\par


\begin{proof}
For each $i$, apply \Cref{lem:ok-mvob-i-model-approximation} and write
$u_i=P_{\theta_i}[\varphi_i]_{\cI}$.  The replacement preserves every
body, and by definition \eqref{eq:ok-mvob-mixed-mass} the left side is the
mixed non-pluripolar mass of the $u_i$.  Choose the approximants
$u_{i,j}$ in \eqref{eq:ok-mvob-approximating}, using one common index for
the finitely many $i$.  At every $j$,
\Cref{prop:ok-mvob-analytic-case} gives
\[
 \int_X\bigwedge_i(\theta_i+\ddc u_{i,j})
 =n!\sup_\nu V(\Delta_\nu(\theta_1,u_{1,j}),\ldots,
                    \Delta_\nu(\theta_n,u_{n,j})).
\]
The mixed non-pluripolar masses on the left converge by
\cite[Theorem~6.2.1]{XiaBook}.  The right side converges by
\Cref{prop:ok-mvob-uniform-ds}; explicitly, for the corresponding functions
$f_j,f$ of $\nu$,
\[
 \left|\sup_\nu f_j(\nu)-\sup_\nu f(\nu)\right|
 \leq\sup_\nu|f_j(\nu)-f(\nu)|\longrightarrow0.
\]
Taking the limit proves \eqref{eq:ok-mvob-metrics}.
\end{proof}

\begin{theorem}\label{thm:ok-mvob-line-bundles}
For big line bundles $L_1,\ldots,L_n$ on $X$,
\[
 \langle L_1,\dots,L_n\rangle
 =n!\,\sup_\nu V\!\left(
     \Delta_\nu(L_1),\dots,\Delta_\nu(L_n)\right).
\tag{55.49}\label{eq:ok-mvob-line-bundles}
\]
\end{theorem}

\begin{proof}
Choose psh metrics with minimal singularities in the classes $c_1(L_i)$.
Their partial bodies are the ordinary Okounkov bodies by
\Cref{thm:ok-algebraic-volume-order}.  The mixed non-pluripolar mass of
minimal currents is the positive intersection number by
\cite[Definition~1.17]{BEGZ10}; its projective numerical interpretation is
the movable product of \cite[Propositions~2.12--2.13]{BFJ09}.
Thus \Cref{thm:ok-mvob-metrics} specializes exactly to
\eqref{eq:ok-mvob-line-bundles}.
\end{proof}

\begin{remark}[The supremum is a supremum]
\label{rem:ok-mvob-supremum}
The proofs assert no maximizing valuation.  In the analytic lower bound the
flag may vary with $\varepsilon$, and the uniformity argument uses only
approximate maximizers.  When all
$L_i=L$ coincide, the supremum is superfluous, since
$V(\Delta_\nu(L),\dots,\Delta_\nu(L))=\lambda_n(\Delta_\nu(L))
 =\vol(L)/n!$ for every $\nu$, by
\Cref{thm:appD-minkowski-polynomial} and by
\eqref{eq:ok-algebraic-volume} applied to a potential with minimal
singularities.  With genuinely different $L_i$ the value
does depend on $\nu$, which is why the supremum appears.
\end{remark}

\par\noindent\begin{minipage}{\linewidth}
\begin{example}\label{ex:ok-mvob-diagonal}
Let $L_1=\cdots=L_n=L$ and $h_1=\cdots=h_n=h$, with positive ideal
volume.  For every valuation in \eqref{eq:ok-mvob-valuations},
\[
 n!V(\Delta_\nu(L,h),\ldots,\Delta_\nu(L,h))
 =n!\lambda_n(\Delta_\nu(L,h))
 =\vol(\ddc h,\ldots,\ddc h).
\]
Thus every valuation computes the common value; the supremum is needed only
for genuinely mixed data.
\end{example}
\end{minipage}\par


\begin{example}\label{ex:ok-mvob-toric-boxes}
On $X=(\mathbf P^1)^n$, let
$L_i=\mathcal O(a_{i1},\ldots,a_{in})$ with all $a_{ir}>0$.  The
standard toric flag gives
$\Delta_\nu(L_i)=\prod_r[0,a_{ir}]$.  The polarization formula for boxes
therefore gives
\[
 n!V(\Delta_\nu(L_1),\ldots,\Delta_\nu(L_n))
 =\operatorname{perm}(a_{ir})=(L_1\cdots L_n).
\]
In this case one fixed toric valuation already realizes the supremum in
\eqref{eq:ok-mvob-line-bundles}.
\end{example}

\begin{warning}{Scope of the theorem}
\label{warn:ok-mixed-scope}
\index{Okounkov body!zero-volume regime}
The preceding proof carries a strict positivity hypothesis and says nothing
outside it.

\emph{Zero volume is excluded.}  The metric claim requires
$\vol(\theta_i,\varphi_i)>0$ for every $i$, and
the line-bundle claim requires every $L_i$ big.  If one factor
has zero volume, the corresponding partial body is not defined by
\Cref{def:ok-algebraic-partial-body}, both sides of
\eqref{eq:ok-mvob-metrics} leave the stated domain, and no continuity
argument in this chapter recovers an identity.  That regime is exactly what
\Cref{prob:frontier-zero-volume} keeps open.

\emph{Transcendental classes are not covered.}  The theorems concern $n$
line bundles on a smooth projective $X$, not transcendental big
classes on a compact K\"ahler manifold.  The transcendental bodies of
\eqref{eq:ok-trans-general} and \Cref{thm:ok-transcendental-package} are
available, and \eqref{eq:ok-minkowski} is their superadditivity, but no
transcendental analogue of \eqref{eq:ok-mvob-metrics} is claimed here.

\emph{Attainment is not asserted.}  See \Cref{rem:ok-mvob-supremum}.
\end{warning}

\begin{remark}[Status]
\label{rem:ok-mixed-status}
\index{Okounkov body!mixed volume!status of Conjecture 8.8}
The mixed partial-body volume formula was posed as
\cite[Conjecture~8.8]{Xia21}.  The author-supplied July 2026 proof
\cite{XiaMVOB}, after the repairs made explicit above, proves it on its
positive-ideal-volume projective domain.  Its scope means that it is not a
zero-volume theorem.  The numerical-dimension question
\cite[Conjecture~5.29]{Xia21} concerns limit partial bodies, including
zero volume.  The Chebyshev-transform question
\cite[Section~6.4, p.~1325]{Xia21} instead assumes positive ideal volume.
Both remain separately discussed in \Cref{prob:frontier-zero-volume}.
\end{remark}

\section{The zero-volume boundary}
\label{sec:ok-zero-volume}

The full-dimensional volume identity cannot simply be continued to
multiplier-ideal volume zero.
The toric convex dictionary nevertheless retains lower-dimensional data.
We first treat arbitrary real classes analytically and then retain the
section-counting argument in integral and rational classes as a second
proof.  The numerical dimensions below belong to the current, not merely to
its cohomology class.

Let $X=X_\Sigma$ be a smooth projective toric variety, and let $M$ and
$N$ be its character and cocharacter lattices.  Invariance in this section
means invariance under the maximal compact torus.  Normalize Lebesgue measure
on $M_{\mathbf R}$ so that a fundamental lattice parallelepiped has volume
one.  Use the logarithmic coordinates characterized by
$\langle u,\operatorname{Log}z\rangle=\log|\chi^u(z)|$, as in the proof of
\Cref{thm:ok-toric-zero-mass}.  If $T$ is an invariant closed positive
$(1,1)$-current, its restriction to the dense torus has the form
\[
             T|_{(\mathbf C^*)^n}=\ddc\bigl(2\,F_T\circ\operatorname{Log}\bigr)
\]
after an invariant smooth local primitive has been absorbed into a finite
convex function $F_T:N_{\mathbf R}\to\mathbf R$, the convex weight of the
\emph{norm}; the squared-norm weight of \Cref{sec:ch51-squared-norm} is
$2F_T\circ\operatorname{Log}$.  The factor $2$ is what makes
$\Delta(\omega_A)$ the moment polytope of $A$ rather than twice it, and
hence what makes \eqref{eq:ok-toric-real-products} hold with no factor
$2^n$: on $\mathbf P^1$ with $T=\omega$, dropping it would give
$K=[0,2]$ and turn the identity $1=\int_X\langle T\rangle=1!\,V(K)$ into
$1=2$.  Changing the primitive
adds an affine function to $F_T$.

\begin{definition}\label{def:ok-closed-newton-body}
The closed Newton body of $T$ is
\[
                 \Delta(T):=
                 \overline{\operatorname{dom}F_T^*}
                 \subset M_{\mathbf R}.
\tag{55.50}\label{eq:ok-closed-newton-body}
\]
It is understood up to translation.
\end{definition}

The closure is load bearing.  For example,
$F(x)=\log(1+\rme^x)$ has gradient image $(0,1)$, whereas
$\operatorname{dom}F^*=[0,1]$.  Dimension and mixed volume see the closed
body.  This definition agrees with the usual toric Newton body on the
positive-volume locus.  At zero volume it is a separate toric object and is
not being defined through \Cref{def:ok-algebraic-partial-body}.

For a current $T$, put
\[
 \nu_{\mathrm{Cao}}(T)
 :=\max\{k:\langle T^k\rangle_{\mathrm{Cao}}\ne0\},
 \qquad
 \nu_{\mathrm{np}}(T)
 :=\max\{k:[\langle T^k\rangle_{\mathrm{np}}]\ne0\}.
\]
The first product is Cao's current-level cohomological product from
quasi-equisingular approximations \cite[Definitions~2--3]{Cao14}; the second
is the ordinary non-pluripolar product.  We use
$\langle T^0\rangle_{\mathrm{Cao}}
=\langle T^0\rangle_{\mathrm{np}}=1$.

\par\noindent\begin{minipage}{\linewidth}
\begin{theorem}\label{thm:ok-toric-real-products}
Let $X=X_\Sigma$ be a smooth projective toric variety of dimension
$n\geq1$, and let $T$ be an invariant closed positive
$(1,1)$-current in a pseudoeffective real class.  Put
$K=\Delta(T)$.  Then $K$ is nonempty and compact.

Let $\omega$ be an invariant K\"ahler form with moment polytope $Q$,
and normalize mixed volume by
$V(B[n])=\operatorname{vol}_M(B)$.  For every $0\leq k\leq n$,
\[
 \int_X\langle T^k\rangle_{\mathrm{Cao}}\wedge\omega^{n-k}
 =\int_X\langle T^k\rangle_{\mathrm{np}}\wedge\omega^{n-k}
 =n!V\bigl(K[k],Q[n-k]\bigr).
\tag{55.51}\label{eq:ok-toric-real-products}
\]
Consequently,
\[
                    \nu_{\mathrm{Cao}}(T)
                    =\nu_{\mathrm{np}}(T)
                    =\dim K.
\tag{55.52}\label{eq:ok-toric-real-dimension}
\]
\end{theorem}
\end{minipage}\par


\begin{proof}
Choose an invariant smooth representative $\theta$ of $\{T\}$, average
a potential over the compact torus, and write
\[
                         T=\theta+\ddc\varphi.
\]
On the dense torus, an invariant primitive of $\theta$, together with
$\varphi$, gives the finite convex function $F=F_T$.  A finite toric
affine cover and the extension of $T$ across its boundary give
$F(x)\leq C(1+\lVert x\rVert)$.  A subgradient exists at every point
because $F$ is finite and convex on all of $N_{\mathbf R}$.  The linear
upper bound makes $\operatorname{dom}F^*$ bounded.  Hence
\[
 K=\overline{\operatorname{dom}F^*}
  =\overline{\bigcup_{x\in N_{\mathbf R}}\partial F(x)}
\]
is nonempty and compact.  Adding an affine function to $F$ translates
this body, proving the asserted intrinsicness.

Fix a decreasing quasi-equisingular approximation $(\varphi_j)$ with
gentle analytic singularities and put
\[
                       R_j=\theta+\ddc\varphi_j.
\]
Such an approximation exists by
\cite[Theorem~1.6.2 and Definition~1.6.5]{XiaBook}.  After comparing its
auxiliary K\"ahler form with $\omega$, we may write
$R_j\geq-\varepsilon_j\omega$, where $\varepsilon_j\downarrow0$.
The multiplier-ideal convention in \cite{Cao14} differs from that of
\cite{XiaBook} by a factor of two.  More precisely,
\[
             \mathcal I_{\mathrm C}(a\varphi)
             =\mathcal I_{\mathrm B}(2a\varphi).
\]
Applying the book's quasi-equisingular inclusion with
$2m<2m(1+\delta)$, and using that the approximants decrease, shows that
the same tail is admissible in Cao's definition.

For $s>0$, set $S_s=T+s\omega$.  Deleting finitely many terms makes
$(R_j+s\omega)$ a positive quasi-equisingular approximation of $S_s$.
Addition of a smooth form changes neither analytic singularities nor
multiplier ideals, and
\[
              (R_j+s\omega)_{\mathrm{ac}}
              =(R_j)_{\mathrm{ac}}+s\omega.
\]
Cao's definition and the binomial formula therefore give
\[
 \int_X\langle S_s^n\rangle_{\mathrm{Cao}}
 =\sum_{k=0}^n\binom nk s^{n-k}
   \int_X\langle T^k\rangle_{\mathrm{Cao}}
                  \wedge\omega^{n-k}.
\tag{55.53}\label{eq:ok-cao-smooth-polynomial}
\]
This is an identity for one fixed approximation, not a continuity assertion
under variation of the cohomology class.

We next identify its left-hand side.  The current $S_s$ dominates
$s\omega$, so it is a K\"ahler current.  Its potential is toric and hence
$\mathcal I$-good by \cite[Example~7.4.1]{XiaBook}.  Thus
\cite[Theorem~7.1.1]{XiaBook} gives $d_{\mathcal S}$-convergence of the
fixed quasi-equisingular approximation, and
\cite[Theorem~6.2.1]{XiaBook} gives convergence of its top non-pluripolar
masses.

Gentleness is essential at this point.  For fixed $s$ and every term in
the positive tail, there is a modification $\pi_j:Y_j\to X$ and an exact
decomposition
\[
                    \pi_j^*(R_j+s\omega)
                    =[E_j]+\beta_{j,s},
\]
where $E_j$ is an effective real divisor and $\beta_{j,s}$ is smooth.
Off $\operatorname{Supp}E_j$, the form $\beta_{j,s}$ is the pullback of
a positive current, so it is semipositive there and, by continuity,
everywhere.  The singular and exceptional analytic sets have Lebesgue
measure zero, while the non-pluripolar product ignores them.  On their
complement both products are the ordinary smooth wedge.  It follows that
\[
 \int_X(R_j+s\omega)_{\mathrm{ac}}^n
 =\int_{Y_j}\beta_{j,s}^n
 =\int_X\langle(R_j+s\omega)^n\rangle_{\mathrm{np}}.
\]
Taking the limit gives
\[
 \int_X\langle S_s^n\rangle_{\mathrm{Cao}}
 =\int_X\langle S_s^n\rangle_{\mathrm{np}}.
\tag{55.54}\label{eq:ok-cao-np-positive-perturbation}
\]
A merely bounded remainder in a general analytic singularity would not
provide the smooth residual form used above.

Let $G$ be the smooth strictly convex weight of $\omega$, so that
$\overline{\operatorname{dom}G^*}=Q$.  The weight of $S_s$ is
$F_s=F+sG$, and
\[
             \overline{\operatorname{dom}F_s^*}=K+sQ.
\tag{55.55}\label{eq:ok-real-newton-minkowski}
\]
For completeness, if $H$ is a finite convex function with bounded closed
conjugate domain $B$, its recession function is the support function
$H^\infty=h_B$.  Recession functions add, so
\[
                 F_s^\infty=F^\infty+sG^\infty
                 =h_K+s h_Q=h_{K+sQ}.
\]
Support functions determine compact convex bodies, proving
\eqref{eq:ok-real-newton-minkowski}.  This argument avoids the generally
false unclosed-domain identity for conjugates of sums.

The function $F_s$ is locally bounded on $N_{\mathbf R}$.  The local
complex/real Monge--Amp\`ere calculation
\cite[Lemmas~2.1--2.2]{CGSZ19} identifies the logarithmic pushforward of its
Bedford--Taylor measure with $n!$ times its Alexandrov measure.  The
interior of $\operatorname{dom}F_s^*$ lies in the subgradient image, and
the latter lies in the domain.  Since $K+sQ$ is full dimensional, closing
that image changes only a Lebesgue-null boundary.  The toric boundary of
$X$ is analytic and hence pluripolar, so the non-pluripolar measure does
not charge it \cite[Propositions~1.4 and~1.6]{BEGZ10}.  Consequently,
\[
 \int_X\langle S_s^n\rangle_{\mathrm{np}}
       =n!\operatorname{vol}_M(K+sQ).
\tag{55.56}\label{eq:ok-real-toric-mass}
\]
Thus only the local torus calculation has been used; no global
ample-class theorem has been extended silently to a real class.

Minkowski's polynomial formula reads
\[
 \operatorname{vol}_M(K+sQ)
 =\sum_{k=0}^n\binom nk
       V(K[k],Q[n-k])s^{n-k}.
\tag{55.57}\label{eq:ok-real-minkowski-polynomial}
\]
Equations \eqref{eq:ok-cao-smooth-polynomial},
\eqref{eq:ok-cao-np-positive-perturbation},
\eqref{eq:ok-real-toric-mass}, and
\eqref{eq:ok-real-minkowski-polynomial} hold for every $s>0$.  Equality of
polynomials gives
\[
 \int_X\langle T^k\rangle_{\mathrm{Cao}}\wedge\omega^{n-k}
       =n!V(K[k],Q[n-k]).
\]
Multilinearity of non-pluripolar products under addition of a smooth form
\cite[Proposition~1.4]{BEGZ10} gives, by the same coefficient comparison,
\[
 \int_X\langle T^k\rangle_{\mathrm{np}}\wedge\omega^{n-k}
       =n!V(K[k],Q[n-k]).
\]
This proves \eqref{eq:ok-toric-real-products}.

It remains to determine which coefficients vanish.  Put $d=\dim K$.  If
$k>d$, translate $K$ into a $d$-dimensional linear subspace.  The
volume polynomial has degree at most $d$ in the variables belonging to
the $k$ copies of $K$, and hence
$V(K[k],Q[n-k])=0$.  If $1\leq k\leq d$, choose $k$ linearly
independent segments in a translate of $K$, and complete their directions
by $n-k$ segments contained, after translation and scaling, in the
interior of $Q$.  The mixed volume of the resulting segments is a positive
multiple of the absolute determinant of their directions.  Monotonicity of
mixed volume therefore gives
\[
 V(K[k],Q[n-k])
 \begin{cases}
  >0,&0\leq k\leq d,\\
  =0,&d<k\leq n.
 \end{cases}
\tag{55.58}\label{eq:ok-real-mixed-positivity}
\]
Cao's product class is pseudoeffective, and the non-pluripolar product is
represented by a positive current.  Either class is nonzero exactly when
its pairing with the strictly positive form $\omega^{n-k}$ is positive.
The two coefficient formulas and
\eqref{eq:ok-real-mixed-positivity} prove
\eqref{eq:ok-toric-real-dimension}.
\end{proof}

\begin{example}\label{ex:ok-real-toric-class-current}
On $X=\mathbf P^1$, let $T=\lambda[0]$, with $\lambda>0$.  The class
$\lambda[\omega_{\mathrm{FS}}]$ has class numerical dimension one, but
the Newton body of $T$ is a point and its current-level numerical
dimensions are zero.  Taking $\lambda\notin\mathbf Q$ gives the same
counterexample in an irrational class.  Thus
\Cref{thm:ok-toric-real-products} would be false if
$\nu(T)$ were interpreted as the numerical dimension of $\{T\}$.
\end{example}

\begin{example}\label{ex:ok-real-toric-segment}
On $\mathbf P^1\times\mathbf P^1$, take
\[
 T=a[\{0\}\times\mathbf P^1]
      +b\,\operatorname{pr}_2^*\omega_{\mathrm{FS}},
 \qquad a,b>0,\quad a/b\notin\mathbf Q.
\]
Up to translation, $\Delta(T)=\{0\}\times[0,b]$.  For the product
K\"ahler form with square moment polytope,
$\operatorname{vol}(\Delta(T)+sQ)=bs+s^2$.  Hence both current-level
numerical dimensions are one, while the big cohomology class has numerical
dimension two.
\end{example}

The next rational statement is already contained in
\Cref{thm:ok-toric-real-products}.  Its proof is retained because it gives an
independent algebraic explanation through multiplier ideals and character
counts.

\par\noindent\begin{minipage}{\linewidth}
\begin{corollary}\label{thm:ok-toric-zero-mass}
Let $X$ be a smooth projective toric variety.
Let $T$ be a torus-invariant positive $(1,1)$-current in an integral or
rational big toric class, and let
\[
                             K=\Delta(T)
\]
be its Newton body, in the separate toric zero-volume sense described
above.  No use of \Cref{def:ok-algebraic-partial-body} at zero volume is
intended.
Then Cao's numerical dimension satisfies
\[
                         \nu_{\mathrm{Cao}}(T)=\dim K.
\tag{55.59}\label{eq:ok-toric-numdim}
\]
\end{corollary}
\end{minipage}\par


\begin{proof}
Choose a toric $\mathbf Q$-divisor $D$ representing the class of $T$.
After multiplication by a common denominator, we may suppose that $D$ is
integral and that $T$ is the curvature current of a torus-invariant
semipositive singular metric $h_0\rme^{-\varphi}$ on
$\mathcal O_X(D)$.  This operation replaces $K$ by a positive homothetic
copy and does not change either $\dim K$ or Cao's numerical dimension.  Let
$M$ and $N$ be the character and cocharacter lattices.  We normalize
Lebesgue measure $\lambda_M$ on $M_{\mathbf R}$ so that a fundamental
parallelepiped of $M$ has measure one.

Choose a very ample integral toric divisor $A$, with full-dimensional
moment polytope $Q$, and a torus-invariant K\"ahler form $\omega_A$ in
its class.  We use logarithmic coordinates characterized by
$\langle u,\operatorname{Log}z\rangle=\log|\chi^u(z)|$, and let
$F_T$ and $F_A$
be the corresponding convex weights of the \emph{norms} on
$N_{\mathbf R}$.  Thus a local squared-norm weight in the convention of
\Cref{sec:ch51-squared-norm} is twice the corresponding norm weight.  The
Newton body is the closure of the effective domain of the convex conjugate
of $F_T$.  The recession-function argument in
\eqref{eq:ok-real-newton-minkowski} gives the closed-domain Minkowski
identity.  Consequently, for rational $\varepsilon>0$,
\[
                  \Delta(T+\varepsilon\omega_A)=K+\varepsilon Q.
\tag{55.60}\label{eq:ok-toric-minkowski}
\]

We now prove the section-counting identity needed for Cao's criterion.  Put
$B_\varepsilon=K+\varepsilon Q$, and take $m$ sufficiently divisible that
$m\varepsilon\in\mathbf Z$.  Set
\[
 W_{m,\varepsilon}:=
 H^0\!\left(X,\mathcal O_X(mD+m\varepsilon A)
                 \otimes\mathcal I(m\varphi)\right).
\]
The multiplier ideal is torus invariant, so $W_{m,\varepsilon}$ is the
direct sum of its one-dimensional character spaces.  We determine their
weights chart by chart.

Let $\sigma$ be a maximal cone of the smooth fan of $X$.  Its primitive
ray generators form a basis of $N$; let
$e_{\sigma,1},\ldots,e_{\sigma,n}$ be the dual basis of $M$, and put
$\kappa_\sigma=\sum_j e_{\sigma,j}$.  Let
$d_\sigma,a_\sigma\in M$ be the Cartier data of $D,A$ on the affine chart
$U_\sigma$.  In the corresponding toric frame, the coefficient of the
character section $\chi^u$ of
$\mathcal O_X(mD+m\varepsilon A)$ is the monomial with exponent
\[
       \alpha_\sigma:=u-m d_\sigma-m\varepsilon a_\sigma.
\]
Let $\Phi_\sigma$ be the local squared-norm psh weight of $T$ in this
frame, and put $\psi_\sigma=\Phi_\sigma/2$.  Thus $\psi_\sigma$ is the
norm weight used by Guenancia.  The global relative potential $\varphi$
and $\Phi_\sigma$ differ by a smooth function, and hence
\[
 \mathcal I(m\varphi)|_{U_\sigma}
   =\mathcal I(m\Phi_\sigma)
   =\mathcal I_{\mathrm G}(m\psi_\sigma),
\]
where the last ideal is defined by local integrability against
$\rme^{-2m\psi_\sigma}$.  The closure of the local Newton convex set of
$\psi_\sigma$ is
\[
              \overline{P(\psi_\sigma)}
                    =K-d_\sigma+\sigma^\vee.
\]
Indeed, subtracting the Cartier linear function translates the effective
domain by $-d_\sigma$, while allowing holomorphic monomials adds the dual
cone.  The set $P(\psi_\sigma)$ is convex and the displayed closure is
full-dimensional.  A convex set with full-dimensional closure has the same
interior as its closure, and hence
\[
 \operatorname{int}P(m\psi_\sigma)
   =\operatorname{int}\bigl(m(K-d_\sigma)+\sigma^\vee\bigr).
\]

Guenancia's toric multiplier-ideal criterion
\cite[Theorem~1.13]{Guenancia12} says that a holomorphic monomial of exponent
$\alpha_\sigma$ belongs to $\mathcal I_{\mathrm G}(m\psi_\sigma)$
exactly when $\alpha_\sigma+\kappa_\sigma$ lies in this interior.  This
condition also recovers holomorphicity, so it may be applied to an arbitrary
character $\chi^u$.  Indeed, if
\[
 P_D:=\bigcap_{\tau\in\Sigma(n)}(d_\tau+\tau^\vee)
\]
is the toric polytope of $D$, extension of the invariant psh metric across
the toric charts gives $K\subseteq P_D$.  Hence
$m(K-d_\sigma)+\sigma^\vee\subseteq\sigma^\vee$.  Membership of
$\alpha_\sigma+\kappa_\sigma$ in its interior implies, in the smooth chart
basis, $\alpha_{\sigma,j}+1>0$ for every $j$.  Since these coordinates
are integers, $\alpha_{\sigma,j}\geq0$, as required.

After cancelling $m d_\sigma$, we obtain
\[
 \chi^u\in W_{m,\varepsilon}
 \quad\Longleftrightarrow\quad
 u+\kappa_\sigma\in
 \operatorname{int}(mB_\varepsilon+\sigma^\vee)
 \quad\text{for every maximal }\sigma.
\]
Here we used
$Q+\sigma^\vee=a_\sigma+\sigma^\vee$, which holds because $A$ is
very ample.  Notice that the criterion includes both the strict interior and
the fixed anticanonical shift $\kappa_\sigma$; neither may be discarded at
finite level.

Let $E_{m,\varepsilon}\subset M$ be the set on the right.  Then
$\dim W_{m,\varepsilon}=\#E_{m,\varepsilon}$ by the toric character
dictionary \cite[\S4.3]{CLS11}.  Completeness of the fan gives the elementary
gluing identity
\[
 B_\varepsilon
   =\bigcap_{\sigma\in\Sigma(n)}
        (B_\varepsilon+\sigma^\vee).
\]
To see the reverse inclusion, if $x\notin B_\varepsilon$, strict separation
gives $v\in N_{\mathbf R}$ such that
$\langle x,v\rangle<\min_{b\in B_\varepsilon}\langle b,v\rangle$.
Choose a maximal cone containing $v$.  Every point of
$B_\varepsilon+\sigma^\vee$ has pairing with $v$ at least that minimum,
so $x$ is excluded from that member of the intersection.

After division by $m$, the exact character conditions cut out
\[
 C_{m,\varepsilon}:=
 \bigcap_{\sigma\in\Sigma(n)}
 \left(-\frac{\kappa_\sigma}{m}
       +\operatorname{int}(B_\varepsilon+\sigma^\vee)\right),
 \qquad
 E_{m,\varepsilon}=mC_{m,\varepsilon}\cap M.
\]
The fan is finite and the shifts are $O(m^{-1})$.  Hence every compact
subset of $\operatorname{int}B_\varepsilon$ lies in
$C_{m,\varepsilon}$ for all large $m$, while
$C_{m,\varepsilon}$ lies in any prescribed neighborhood of
$B_\varepsilon$ for all large $m$.  For completeness, the second
assertion follows by contradiction: the intersection is uniformly bounded
because cones containing both signs of each coordinate direction give
uniform upper and lower bounds; a convergent sequence staying a fixed
distance from $B_\varepsilon$ would, after passage to the limit, belong to
every closed set $B_\varepsilon+\sigma^\vee$, hence to
$B_\varepsilon$, a contradiction.

Since $\varepsilon>0$ and $Q$ has interior, $B_\varepsilon$ is
full-dimensional and its boundary has $\lambda_M$-measure zero.  Choose
$b\in\operatorname{int}B_\varepsilon$.  The homothetic inner bodies
$b+(1-\delta)(B_\varepsilon-b)$ and arbitrarily small outer parallel
bodies squeeze $C_{m,\varepsilon}$.  The lattice Riemann-sum formula for a
bounded convex body $C$,
\[
             \frac{\#(mC\cap M)}{m^n}\longrightarrow\lambda_M(C),
\]
therefore yields
\[
 \lim_{\substack{m\to\infty\\m\varepsilon\in\mathbf Z}}
 \frac{\dim W_{m,\varepsilon}}{m^n}
       =\lambda_M(K+\varepsilon Q).
\]
One elementary proof of the Riemann-sum formula surrounds each lattice point
by its covolume-$m^{-n}$ fundamental cell; the symmetric difference is
contained in a shrinking neighborhood of $\partial C$.  This also explains
why no factor $n!$ occurs here.  We have proved an actual limit, not a
finite-level equality or merely a limsup.

Since the preceding inner limit exists, Cao's projective criterion
\cite[Proposition~4.4]{Cao14} becomes, in this normalization,
\[
 \lim_{\substack{\varepsilon\downarrow0\\
                  \varepsilon\in\mathbf Q}}
 \frac{\log\!\left(
   \displaystyle\lim_{\substack{m\to\infty\\m\varepsilon\in\mathbf Z}}
   \frac{h^0\!\left(X,\mathcal O_X(mD+m\varepsilon A)
                    \otimes\mathcal I(m\varphi)\right)}{m^n}
   \right)}{\log\varepsilon}
 =n-\nu_{\mathrm{Cao}}(T).
\]
The ordinary multiplier ideal in this formula is exactly the one counted
above; Cao proves that its asymptotics agree with those of the
upper-regularized ideal used during his proof.

Put $d=\dim K$.
After translation, $K$ contains a $d$-dimensional ball in its affine
span and is contained in a bounded $d$-dimensional box.  The body $Q$
contains an $n$-dimensional ball and is contained in a bounded
$n$-dimensional box.  Splitting $M_{\mathbf R}$ into the direction space
of $K$ and a complement, and comparing with products of balls and boxes,
gives constants $c,C>0$ such that
\[
 c\varepsilon^{\,n-d}
 \leq\lambda_M(K+\varepsilon Q)
 \leq C\varepsilon^{\,n-d}
\tag{55.61}\label{eq:ok-toric-tube-growth}
\]
for all sufficiently small $\varepsilon>0$.
Thus
$\log\lambda_M(K+\varepsilon Q)
 =(n-d)\log\varepsilon+O(1)$.  Substitution in Cao's criterion gives
$n-\nu_{\mathrm{Cao}}(T)=n-d$, hence
$\nu_{\mathrm{Cao}}(T)=d$.
\end{proof}

\begin{remark}[Exact scope]
\label{rem:ok-toric-scope}
The section-counting proof of \Cref{thm:ok-toric-zero-mass} is restricted to
integral or rational big classes because it uses one graded section algebra.
The conclusion itself is not so restricted:
\Cref{thm:ok-toric-real-products} proves it analytically for every invariant
positive current in a pseudoeffective real class.  The two arguments must not
be conflated: the real-class proof fixes the current and adds a smooth form;
it does not approximate the class by rational classes.
\end{remark}

\begin{proposition}\label{prop:ok-i-invariant-obstruction}
There is no assignment
\[
              \varphi\longmapsto B(\varphi)\subset\mathbf R^N
\]
of measurable sets, invariant under $\cI$-equivalence, whose Euclidean
volume equals the actual non-pluripolar mass of every psh potential.
\end{proposition}

\begin{proof}
On $\mathbf P^1$, let $K$ be a polar Cantor set carrying an atom-free
probability measure $\mu$; such a pair is supplied by \cite[p.~31]{Car83}.
Write
\[
                         \mu=\omega+\ddc\varphi
\]
with $\omega$ Fubini--Study.
Because $\mu$ has no atoms, every Lelong number of $\varphi$ is zero.
The valuative formula for multiplier ideals therefore gives
\[
                              \varphi\sim_{\cI}0.
\]
But $\mu$ is supported on the pluripolar set $K$, so its
non-pluripolar mass is zero, whereas
\[
                       \int_{\mathbf P^1}\omega=1.
\]
An $\cI$-invariant assignment would give
$B(\varphi)=B(0)$, while the desired volume rule would give these two sets
different volumes, a contradiction.
\end{proof}

Convexity played no role in the obstruction.
It rules out ordinary multiplier-ideal bodies as universal carriers of
actual mass, but it does not rule out a refinement that retains the finer
$P$-singularity type.
That $P$-sensitive problem remains open.

For the limit partial body defined by ample perturbations,
\cite[Section~5.6 and Conjecture~5.29]{Xia21} predicts equality between
its dimension and the numerical dimension, including at zero volume.
The question in \cite[Section~6.4, p.~1325]{Xia21} has a different domain:
it assumes $\vol(\theta,\varphi)>0$ and replaces $\varphi$ by its
$\mathcal I$-model projection.  Choose $\eta\in\PSH(X,\theta)$ with
$\theta_\eta$ a K\"ahler current and $\eta\preceq\varphi$, and set
$\varphi_\varepsilon=(1-\varepsilon)\varphi+\varepsilon\eta$.
For a fixed continuous metric weight $v$ and flag, the source asks whether
$\lim_{\varepsilon\downarrow0}c_{[\varphi_\varepsilon]}[v]$
on $\operatorname{Int}\Delta_\nu(\theta,\varphi)$ is independent of
the auxiliary potential $\eta$.  It is a positive-volume choice problem.
These source questions and the stronger $P$-sensitive zero-volume
construction are distinguished in \Cref{prob:frontier-zero-volume}.
Conjecture~8.8 of that source is the positive-volume mixed-body formula
proved in \Cref{thm:ok-mvob-metrics}.  This closes the conjecture on its
stated positive-ideal-volume projective domain, but it does not address any
of the zero-volume problems above or the Chebyshev choice problem.

\section{Provenance and safeguards}
\label{sec:ok-provenance}

The algebraic partial-body and test-function theory is from the published
\cite[Sections~3, 5 and~7]{Xia21}.
The transcendental bodies use \cite[Section~10.4]{XiaBook}, including its
ordinary-body interface to \cite{DRWNXZ}.
The positive-mass test-curve and DH package, beyond the published paper's
finite-energy algebraic setting, is \cite[Sections~10.5 and~13.3]{XiaBook}.
The tail law \eqref{eq:ok-intrinsic-tail} proves flag independence directly;
\eqref{eq:ok-dh-weight-normalization} records the factor $n!$ converting
the published body's measure to this book's intersection-mass convention.
The proof of \Cref{thm:ok-mvob-metrics} starts from the author-supplied
argument in \cite{XiaMVOB}.  Its load-bearing external inputs are the
general-valuation partial-body and $d_S$ results of \cite{XiaBook}, the
initial-ideal comparison of \cite{KM19}, the Bernstein--Kushnirenko mixed
degree formula, the nef-approximation property of positive products, the
minimal-current identification of \cite{BEGZ10}, and the common-flag theorem
\cite[Theorem~1.1]{Wil25}.  The proof above supplies the missing
four-inclusion convex lemma, the tail/order-closure argument, the common
finite subsequence, the one-dimensional-factor branch, the multigraded
irrelevant-torsion check, the rational analytic reduction, and the
canonical-section inclusion.  It was verified in production Rethlas against
the exact repaired proof object of 24 August 2026 (SHA-256)
\begin{center}
\small\texttt{9227826d7ba5c641c5fe32cc4daae9acbb3b57274dcc1b356ba8beadd58cad7b}
\end{center}
with verdict correct and no reported gap.
The complete repaired author source (SHA-256)
\begin{center}
\small\texttt{5746eb011417caecfa876c05c6c74fb548e2ae9d086c1441189bbcad02a96594}
\end{center}
then received an independent
source-level verdict correct, again with no reported error or gap.  The
repaired proof object above remains the controlling mathematical receipt.
These are historical receipts for the named frozen proof objects; they are
not independent verification of the current chapter's edited bytes.

\Cref{prop:ok-mvob-superadditivity} holds for every valuation in
\eqref{eq:ok-mvob-valuations} by \cite[Proposition~10.3.10]{XiaBook}; its
order assertion is \eqref{eq:ok-i-order}.  The convex-geometric normalization
of mixed volume is \Cref{thm:appD-minkowski-polynomial}, and the
multiplier-ideal volume of a metric is
\Cref{def:complex-current-volume}; \Cref{rem:ok-mvob-normalizations} carries
out both conversions explicitly rather than assuming them.
The general toric statement \Cref{thm:ok-toric-real-products} is controlled
by production Rethlas run
\nolinkurl{20260825T011636Z_99edb45c54cb}; its source-aware frozen proof
received verdict \texttt{correct}, with no reported error or gap, and an
independent hostile audit found no defect.  The rational corollary
\Cref{thm:ok-toric-zero-mass} retains the separately validated
section-counting proof through Cao's criterion \cite{Cao14}.  Neither
verification is a novelty or priority claim.
The Cantor-set example uses \cite{Car83}.


The affine Newton domain in
\Cref{sec:appH-toric-multiplier-pairs} supplies a related local
convex dictionary. Seo's \cite{Seo26ToricMultiplierV1} describes
adapted pair multiplier modules by character-lattice points whose
canonical translates lie in the strict interior of a logarithmic
Newton domain. Recession inequalities distinguish its closure from
its interior, and exponential tail integrability explains the strict
threshold. This domain can be unbounded and need not be closed;
it is different from a compact partial Okounkov body associated to
a global class and flag. The local lattice criterion therefore gives
no zero-volume extension of the mixed-volume identity or a formula
for the actual nonpluripolar mass. Its normalization and singular
pair boundaries are recorded in
\Cref{sec:provenance-toric-analytic-multiplier}.

Four boundaries should be retained.
The positive-volume continuity theorem is not a zero-volume definition.
The valuative characterization retains its closure.
The $\cI$-obstruction is not an impossibility theorem for
$P$-sensitive refinements.
And the mixed-volume identity is an identity on the positive-volume locus
only: it says nothing at zero volume, which remains
\Cref{prob:frontier-zero-volume}.  No argument in this dependency chain uses
a Deligne-pairing Riemann--Roch formula.
